\documentclass[../../main.tex]{subfiles}% 注意这里的文件路径不能用 ./main.tex ,否则用latexmk编译子文件会报错
\graphicspath{{\subfix{./image/}}{\subfix{../../image/}}} % 指定图片引用路径,后续可以直接使用图片文件名
% 如果图片存放在上级目录,则使用 ../../image/ 作为备用路径

% 例如:
% \begin{figure}[H]
% \centering
% \includegraphics[width=0.3\textwidth]{图.png}
% \caption{}
% \label{figure:图}
% \end{figure}
% 注意:上述\label{}一定要放在\caption{}之后,否则引用图片序号会只会显示??.

\begin{document}

\section{纲与开映射定理}

\subsection{纲与纲推理}

\begin{definition}
设$(\mathscr{X},\rho)$是一个度量空间, 集合$E\subseteq\mathscr{X}$, 称$E$是\textbf{疏的}或\textbf{疏集}或\textbf{疏朗集}, 如果$\overline{E}$的内点是空的.
\end{definition}

\begin{examples}
在$\mathbb{R}^n$上, 有穷点集是疏集. Cantor集是疏集.
\end{examples}

\begin{proposition}
设$(\mathscr{X},\rho)$是一度量空间. $E\subseteq\mathscr{X}$是疏集当且仅当: 对$\forall B(x_0,r_0),\exists B(x_1,r_1)\subseteq B(x_0,r_0)$, 使得$\overline{E}\cap\overline{B}(x_1,r_1)=\varnothing$.
\end{proposition}
\begin{proof}

{\heiti 必要性:} 因为$\overline{E}$无内点, 所以$\overline{E}$不能包含任一球$B(x_0,r_0)$. 从而$\exists x_1\in B(x_0,r_0)$, 使得$x_1\notin\overline{E}$. 又由$\overline{E}$闭, 所以$\overline{E}^c$是开集，进而$\exists\varepsilon_1>0$, 使得$\overline{B}(x_1,\varepsilon_1)\cap\overline{E}=\varnothing$. 取
\begin{align*}
0<r_1<\min(\varepsilon_1,r_0-\rho(x_0,x_1)),
\end{align*}
便有$B(x_1,r_1)\subseteq B(x_0,r_0)$, $\overline{B}(x_1,r_1)\cap\overline{E}=\varnothing$.

{\heiti 充分性:} 若$E$不疏, 即$\overline{E}$有内点, 则$\exists B(x_0,r_0)\subseteq\overline{E}$. 但由假设$\exists B(x_1,r_1)\subseteq B(x_0,r_0)$, 使得$\overline{B}(x_1,r_1)\cap\overline{E}=\varnothing$. 一方面有$B(x_1,r_1)\cap\overline{E}=B(x_1,r_1)$; 另一方面有$B(x_1,r_1)\cap\overline{E}=\varnothing$. 即得矛盾!
\end{proof}

\begin{definition}
在度量空间$(\mathscr{X},\rho)$上, 集合$E$称为\textbf{第一纲的}, 如果$E=\bigcup\limits_{n=1}^{\infty}E_n$, 其中$E_n$是疏集. 不是第一纲的集合称为\textbf{第二纲集}.
\end{definition}

\begin{examples}
在$\mathbb{R}$上, 有理点集是第一纲集. 更一般地, 可数点集总是第一纲集.
\end{examples}

\begin{theorem}[Baire纲定理]\label{theorem:泛函分析--Baire纲定理}
设$(\mathscr{X},\rho)$是完备度量空间，则下列命题成立且等价：
\begin{enumerate}[(1)]
\item $\mathscr{X}$是第二纲集，即不能表示为可数个疏集的并.
\item 若$\{F_n\}$是一列无处稠密的闭集，则$\bigcup\limits_{n=1}^{\infty}F_n$的内部为空.
\item 若$\{G_n\}$是一列在$\mathscr{X}$中稠密的开集，则$\bigcap\limits_{n=1}^{\infty}G_n$也在$\mathscr{X}$中稠密.
\item 若$\{E_n\}$是一列闭集且$\mathscr{X}=\bigcup\limits_{n=1}^{\infty}E_n$，则存在$n_0$使得$E_{n_0}$有非空内部,即$E_{n_0}$有内点.
\end{enumerate}
\end{theorem}
\begin{proof}

首先证明命题(1)成立.
用反证法. 倘若$\mathscr{X}$是第一纲集，即存在一列疏集$\{E_n\}$，使得
\begin{align}
\mathscr{X}=\bigcup\limits_{n=1}^{\infty}E_n. \label{eq:baire1}
\end{align}
对任意的球$B(x_0,r_0)$，存在球$B(x_1,r_1)\subseteq B(x_0,r_0)$且$r_1<1$，使得$\overline{B}(x_1,r_1)\cap\overline{E_1}=\varnothing$；
对球$B(x_1,r_1)$，存在$B(x_2,r_2)\subseteq B(x_1,r_1)$且$r_2<\frac{1}{2}$，使得$\overline{B}(x_2,r_2)\cap(\overline{E_1}\cup\overline{E_2})=\varnothing$；
如此继续下去，对球$B(x_{n-1},r_{n-1})$，存在$B(x_n,r_n)\subseteq B(x_{n-1},r_{n-1})$且$r_n<\frac{1}{n}$，使得$\overline{B}(x_n,r_n)\cap\overline{E_n}=\varnothing$，从而
\begin{align}
\overline{B}(x_n,r_n)\cap\Bigl(\bigcup\limits_{i=1}^{n}\overline{E_i}\Bigr)=\varnothing \quad (\forall n\in\mathbb{N}). \label{eq:baire2}
\end{align}
于是我们得到闭球套
\begin{align*}
\overline{B}(x_0,r_0)\supseteq\overline{B}(x_1,r_1)\supseteq\cdots\supseteq\overline{B}(x_n,r_n)\supseteq\cdots,
\end{align*}
并且
\begin{align}
\rho(x_{n+p},x_n)\leqslant r_n<\frac{1}{n}\quad (\forall n,p\in\mathbb{N}). \label{eq:baire3}
\end{align}
由此可见$\{x_n\}$是基本列，由完备性，存在$x\in\mathscr{X}$使得$\lim\limits_{n\to\infty}x_n=x$. 在\eqref{eq:baire3}中令$p\to\infty$得$\rho(x,x_n)\leqslant r_n$，从而
\begin{align}
x\in\overline{B}(x_n,r_n)\quad (\forall n\in\mathbb{N}). \label{eq:baire4}
\end{align}
联立\eqref{eq:baire2}与\eqref{eq:baire4}可得$x\notin\bigcup\limits_{n=1}^{\infty}\overline{E_n}$，更有$x\notin\bigcup\limits_{n=1}^{\infty}E_n$，这与\eqref{eq:baire1}矛盾. 故$\mathscr{X}$是第二纲集.

接下来证明各命题的等价性.

{\heiti (1) $\Rightarrow$ (3):} 设$\{G_n\}$是一列稠密开集. 倘若$\bigcap\limits_{n=1}^{\infty}G_n$不稠密，则存在开球$U$使得$U\cap\bigl(\bigcap\limits_{n=1}^{\infty}G_n\bigr)=\varnothing$，即$U\subseteq\bigcup\limits_{n=1}^{\infty}(\mathscr{X}\setminus G_n)$. 记$F_n=\mathscr{X}\setminus G_n$，则$F_n$是无处稠密的闭集. 取闭球$B\subseteq U$，则$B$作为完备度量空间的闭子集也是完备的. 且$B=\bigcup\limits_{n=1}^{\infty}(B\cap F_n)$，而$B\cap F_n$在$B$中是无处稠密的闭集. 这表明$B$可表示为可数个疏集的并，即$B$是第一纲集，与(1)矛盾. 故$\bigcap\limits_{n=1}^{\infty}G_n$在$\mathscr{X}$中稠密.

{\heiti (3) $\Rightarrow$ (4):} 设$\{E_n\}$是一列闭集且$\mathscr{X}=\bigcup\limits_{n=1}^{\infty}E_n$. 假设对每个$n$，$E_n$的内部为空，则令$G_n=\mathscr{X}\setminus E_n$，每个$G_n$是开集且在$\mathscr{X}$中稠密. 由(3)知$\bigcap\limits_{n=1}^{\infty}G_n$在$\mathscr{X}$中稠密，但$\bigcap\limits_{n=1}^{\infty}G_n=\mathscr{X}\setminus\bigcup\limits_{n=1}^{\infty}E_n=\varnothing$，矛盾. 故存在某个$E_{n_0}$具有非空内部.

{\heiti (4) $\Rightarrow$ (2):} 设$\{F_n\}$是一列无处稠密的闭集. 若$\bigcup\limits_{n=1}^{\infty}F_n$有非空内部，则存在开球$U\subseteq\bigcup\limits_{n=1}^{\infty}F_n$. 取闭球$B\subseteq U$，则$B$是完备的，且$B=\bigcup\limits_{n=1}^{\infty}(B\cap F_n)$. 每个$B\cap F_n$是$B$中闭集，且因为$F_n$无处稠密，$B\cap F_n$在$B$中也是无处稠密的，从而内部为空. 由(4)应用于完备空间$B$，存在某个$B\cap F_{n_0}$具有非空内部，矛盾. 故$\bigcup\limits_{n=1}^{\infty}F_n$的内部必为空.

{\heiti (2) $\Rightarrow$ (1):} 若$\mathscr{X}$是第一纲集，则存在一列疏集$\{E_n\}$使得$\mathscr{X}=\bigcup\limits_{n=1}^{\infty}E_n$. 令$F_n=\overline{E_n}$，则$F_n$是无处稠密的闭集，且$\mathscr{X}=\bigcup\limits_{n=1}^{\infty}F_n$. 于是$\bigcup\limits_{n=1}^{\infty}F_n$的内部为$\mathscr{X}\neq\varnothing$，与(2)矛盾. 故$\mathscr{X}$是第二纲集.

综上所述，命题(1)(2)(3)(4)彼此等价且都成立.
\end{proof}

\begin{theorem}
在$C[0,1]$中处处不可微的函数集合$E$是非空的, 更确切地, $E$的余集是第一纲集.
\end{theorem}
\begin{remark}
Vander Waerden Function(范德瓦尔登函数):
\begin{align*}
f\left( x \right) =\sum_{n=0}^{\infty}{\frac{\varphi \left( 4^nx \right)}{4^n}},\text{其中}\varphi \left( x \right) =\underset{k\in \mathbb{Z}}{\min}\left\{ \left| x-k \right| \right\}
\end{align*}
就是\href{https://mp.weixin.qq.com/s/ki-x99Me_kKvQUvfXS-GGQ}{处处连续但处处不可微的函数}.
\end{remark}
\begin{proof}
取$\mathscr{X}=C[0,1]$, 设$A_n$表示$\mathscr{X}$中这样一些元素$f$之集: 对$f,\exists s\in[0,1]$, 使对适合$0\leqslant s+h\leqslant1$与$|h|\leqslant1/n$的任何$h$, 成立下式:
\begin{align*}
\left|\frac{f(s+h)-f(s)}{h}\right|\leqslant n.
\end{align*}
若$f$在某个点$s$处可微, 则必有正整数$n$, 使得$f\in A_n$, 于是
\begin{align}
\mathscr{X}\setminus E\subseteq\bigcup\limits_{n=1}^{\infty}A_n.
\end{align}
下面我们证明每个$A_n$是疏集, 为此先证$A_n$是闭的. 事实上, 若$f\in\mathscr{X}\setminus A_n$, 则$\forall s\in[0,1],\exists h_s$, 使得$|h_s|\leqslant\frac{1}{n}$, 且$|f(s+h_s)-f(s)|>n|h_s|$.

又由$f$的连续性, $\exists\varepsilon_s>0$, 以及$s$的某个适当的邻域$J_s$, 使得对$\forall\sigma\in J_s$, 有
\begin{align}
|f(\sigma+h_s)-f(\sigma)|>n|h_s|+2\varepsilon_s.
\end{align}
根据有限覆盖定理, 可设$J_{s_1},J_{s_2},\cdots,J_{s_m}$覆盖$[0,1]$, 并设$\varepsilon=\min\{\varepsilon_{s_1},\varepsilon_{s_2},\cdots,\varepsilon_{s_m}\}$.

今若$g\in\mathscr{X}$适合$\|g-f\|<\varepsilon$, 则由式, 对$\forall\sigma\in J_{s_k}(k=1,2,\cdots,m)$有
\begin{align*}
|g(\sigma+h_{s_k})-g(\sigma)|\geqslant|f(\sigma+h_{s_k})-f(\sigma)|-2\varepsilon>n|h_{s_k}|.
\end{align*}
这证明了$\mathscr{X}\setminus A_n$是开集, 从而$A_n$是闭集.

再证$A_n$没有内点. $\forall f\in A_n,\forall\varepsilon>0$, 由\hyperref[Basis of Analytics-theorem:多项式逼近有限阶光滑函数-Weierstrass逼近定理]{Weierstrass逼近定理}, 存在多项式$p$, 使得$\|f-p\|<\frac{\varepsilon}{2}$.

$p$的导数在$[0,1]$上是有界的, 因此根据中值定理, $\exists M>0$, 使得对$\forall s\in[0,1]$及$|h|<\frac{1}{n}$, 成立$|p(s+h)-p(s)|\leqslant M|h|$.

设$g(s)\in C[0,1]$是一个分段线性函数, 满足$\|g\|<\frac{\varepsilon}{2}$. 并且各条线段斜率的绝对值都大于$M+n$, 那么$p+g\in B(f,\varepsilon)$, 而$p+g\notin A_n$.

这样, 我们证明了每个$A_n$是疏集, 从而$\bigcup\limits_{n=1}^{\infty}A_n$是第一纲集. 而$\mathscr{X}$是完备的, 由\hyperref[theorem:泛函分析--Baire纲定理]{Baire纲定理}$\mathscr{X}$是第二纲集, 由此根据式, $E$也是第二纲集.
\end{proof}

\begin{definition}[Hamel基]
令 \(V\) 是数域 \(\mathbb{F}\) 上的线性空间。如果 \(V\) 的一个子集 \(B \subseteq V\) 满足以下两个条件，则称 \(B\) 为 \(V\) 的一个\textbf{Hamel基}或\textbf{代数基}:
\begin{enumerate}[(1)]
\item \textbf{线性无关性：} 对 \(B\) 中的任意非空有限子集 \(\{b_1, b_2, \cdots, b_n\}\)，若存在标量 \(c_1, c_2, \cdots, c_n \in \mathbb{F}\) 使得 \(\sum\limits_{i=1}^{n} c_i b_i = 0\)，则必有 \(c_1 = c_2 = \cdots = c_n = 0\)；

\item \textbf{完备性：} 对任意向量 \(x \in V\)，都存在 \(B\) 中的一个非空有限子集 \(\{b_1, b_2, \cdots, b_n\}\) 和标量 \(c_1, c_2, \cdots, c_n \in \mathbb{F}\)，使得 \(x\) 可以唯一地表示为有限线性组合 \(x = \sum\limits_{i=1}^{n} c_i b_i\).
\end{enumerate}
\end{definition}

\begin{proposition}
任何一个无穷维 Banach 空间都不存在可数的 Hamel 基.
\end{proposition}
\begin{remark}
显然有限维Banach空间的任意基就是Hamel基.
\end{remark}
\begin{proof}
反设无穷维 Banach 空间 \(X\) 具有可数的 Hamel 基 \(\{e_1, e_2, \cdots\}\).
令 \(X_n = \operatorname{span}\{e_1, e_2, \cdots, e_n\}\)，则对于每一个 \(n \in \mathbb{N}\) 成立. 因为 \(X\) 是无穷维的，所以每个 \(X_n\) 都是 \(X\) 的有限维真子空间. 由\cref{corollary:有限维B*空间必完备子空间必是闭的},在赋范空间中，有限维子空间必为闭子空间.
根据 Hamel 基的定义，\(X\) 中的任意元素都可以表示为 \(\{e_1, e_2, \cdots\}\) 的有限线性组合，故
\begin{align*}
X = \bigcup\limits_{n=1}^{\infty} X_n.
\end{align*}
然而，由\hyperref[theorem:泛函分析--Baire纲定理]{Baire纲定理}可知，一个完备的度量空间不能表示为可数个无处稠密集的并集. 由于每个 \(X_n\) 的内部均为空，故它们都是无处稠密集，这说明 \(X = \bigcup\limits_{n=1}^{\infty} X_n\) 与 \hyperref[theorem:泛函分析--Baire纲定理]{Baire纲定理}相矛盾!
综上，无穷维 Banach 空间不存在可数 Hamel 基，其基必然是无限的且不可数.
\end{proof}

\begin{definition}[Schauder基]
令 \(X\) 是赋范线性空间。如果存在可数序列 \(\{e_n\}_{n=1}^{\infty} \subseteq X\)，使得对任意向量 \(x \in X\)，都存在唯一的标量序列 \(\{c_n\}_{n=1}^{\infty}\)，使得
\begin{align*}
x = \sum\limits_{n=1}^{\infty} c_n e_n,
\end{align*}
这里的级数按范数收敛，即
\begin{align*}
\lim\limits_{N \to \infty} \left\| x - \sum\limits_{n=1}^{N} c_n e_n \right\| = 0,
\end{align*}
则称 \(\{e_n\}_{n=1}^{\infty}\) 为 \(X\) 的一个 \textbf{Schauder基}或\textbf{拓扑基}.
\end{definition}

\begin{proposition}
若赋范线性空间 \(X\) 存在可数 Schauder 基，则 \(X\) 是可分的.
\end{proposition}

\begin{proof}
设 \(\{e_n\}_{n=1}^{\infty}\) 是 \(X\) 的一个可数 Schauder 基. 考虑集合
\begin{align*}
D = \left\{ \sum\limits_{i=1}^{n} q_i e_i : n \in \mathbb{N}, q_i \in \mathbb{Q} \right\}.
\end{align*}
由于有理数集 \(\mathbb{Q}\) 是可数的，\(D\) 是有限个有理数系数的组合，故 \(D\) 是可数集.

下面证明 \(D\) 在 \(X\) 中稠密. 对任意 \(x \in X\) 和任意 \(\varepsilon > 0\)，由于 \(\{e_n\}_{n=1}^{\infty}\) 是 Schauder 基，存在正整数 \(N\)，使得
\begin{align*}
\left\| x - \sum\limits_{i=1}^{N} c_i e_i \right\| < \frac{\varepsilon}{2}.
\end{align*}
由有理数的稠密性，对每个 \(i = 1, 2, \cdots, N\)，存在有理数 \(q_i \in \mathbb{Q}\)，使得
\begin{align*}
|q_i - c_i| < \frac{\varepsilon}{2N \|e_i\|},
\end{align*}
这里不妨设 \(e_i \neq 0\). 于是
\begin{align*}
\left\| x - \sum\limits_{i=1}^{N} q_i e_i \right\| &\leqslant \left\| x - \sum\limits_{i=1}^{N} c_i e_i \right\| + \left\| \sum\limits_{i=1}^{N} (c_i - q_i) e_i \right\| \\
&< \frac{\varepsilon}{2} + \sum\limits_{i=1}^{N} |c_i - q_i| \|e_i\| \\
&< \frac{\varepsilon}{2} + \sum\limits_{i=1}^{N} \frac{\varepsilon}{2N} = \varepsilon.
\end{align*}
因此 \(D\) 在 \(X\) 中稠密. 综上，\(X\) 存在可数稠密子集，故 \(X\) 是可分的.
\end{proof}

\subsection{开映射定理}

\begin{definition}
设$\mathscr{X},\mathscr{Y}$为非空集合，$T:\mathscr{X}\to\mathscr{Y}$为映射。
\begin{enumerate}[(1)]
\item 若存在映射$T_r^{-1}:\mathscr{Y}\to\mathscr{X}$，使得
\begin{align*}
TT_r^{-1} = I,
\end{align*}
其中$I$为$\mathscr{Y}$上的恒同算子，则称$T$\textbf{有右逆}，并称$T_r^{-1}$为$T$的一个\textbf{右逆}。
\item 若存在映射$T_l^{-1}:\mathscr{Y}\to\mathscr{X}$，使得
\begin{align*}
T_l^{-1}T = I,
\end{align*}
其中$I$为$\mathscr{X}$上的恒同算子，则称$T$\textbf{有左逆}，并称$T_l^{-1}$为$T$的一个\textbf{左逆}。
\item 若$T$既有左逆$T_l^{-1}$又有右逆$T_r^{-1}$，则$T_l^{-1}=T_r^{-1}$，此时称$T$是\textbf{可逆的}，并称该公共逆为$T$的\textbf{逆算子}，记为$T^{-1}$，即
\begin{align*}
T_l^{-1} = T_r^{-1} = T^{-1}.
\end{align*}
\end{enumerate}
\end{definition}

\begin{definition}
设$\mathscr{X},\mathscr{Y}$为拓扑空间，映射$T:\mathscr{X}\to\mathscr{Y}$称为\textbf{开映射}，如果$T$将$\mathscr{X}$中的每个开集映为$\mathscr{Y}$中的开集，即对任意开集$U\subseteq\mathscr{X}$，像集$T(U)\subseteq\mathscr{Y}$是开集。
\end{definition}

\begin{definition}
设$\mathscr{X},\mathscr{Y}$都是$B$空间, $T\in\mathscr{L}(\mathscr{X},\mathscr{Y})$.
\begin{enumerate}[(1)]
\item $T$称为是\textbf{单射}, 是指$T$是1-1的.
\item $T$称为是\textbf{满射}, 是指$T(\mathscr{X})=\mathscr{Y}$.
\end{enumerate}
\end{definition}

\begin{theorem}[开映射定理]\label{theorem:泛函分析--开映射定理}
设$\mathscr{X},\mathscr{Y}$都是$B$空间, 若$T\in\mathscr{L}(\mathscr{X},\mathscr{Y})$是一个满射, 则$T$是开映射.
\end{theorem}
\begin{remark}
定理中的$B$空间$\mathscr{X},\mathscr{Y}$可以换成更一般的$F$空间, 但证明需稍做修改. 参看关肇直、张恭庆、冯德兴所著《线性泛函分析入门》(上海科学技术出版社, 1979)的第二章§2.
\end{remark}
\begin{proof}
用$B(x_0,a),U(y_0,b)$分别表示$\mathscr{X},\mathscr{Y}$中的开球.
\begin{enumerate}[(1)]
\item 为了证明$T$是开映射, 即$\forall$开集$W,T(W)$是开集, 当且仅当证明: $\exists\delta>0$, 使得
\begin{align}
TB(\theta,1)\supseteq U(\theta,\delta).
\end{align}
事实上, 必要性是显然的. 下证其充分性. 由于$T$的线性, 充分性假设条件等价于存在$\delta>0$,使得
\begin{align*}
TB(x_0,r)\supseteq U(Tx_0,r\delta)\quad(\forall x_0\in\mathscr{X},\forall r>0).
\end{align*}
$\forall y_0\in T(W)$, 按定义$\exists x_0\in W$, 使得$y_0=Tx_0$. 因为$W$是开集, 所以$\exists B(x_0,r)\subseteq W$. 于是取$\varepsilon=r\delta$, 便有$U(Tx_0,\varepsilon)\subseteq TB(x_0,r)\subseteq T(W)$, 即$y_0=Tx_0$是$T(W)$的内点.
\begin{figure}[H]
\centering
\includegraphics[scale=0.3]{泛函分析图2.3.1.png}
\caption{}
\label{figure:2.3.1}
\end{figure}

\item 证明: $\exists\delta>0$, 使得$\overline{TB(\theta,1)}\supseteq U(\theta,3\delta)$. 这是因为
\begin{align*}
\mathscr{Y}=T\mathscr{X}=\bigcup\limits_{n=1}^{\infty}TB(\theta,n),
\end{align*}
而$\mathscr{Y}$是完备的, 由\hyperref[theorem:泛函分析--Baire纲定理]{Baire纲定理}知$\mathscr{Y}$是第二纲集,所以至少有一个$n\in\mathbb{N}$, 使得$TB(\theta,n)$非疏, 即$\overline{TB(\theta,n)}$至少含有一个内点. 因此$\exists U(y_0,r)\subseteq\overline{TB(\theta,n)}$, 由$T$的线性性质知$TB(\theta,n)$是一个对称凸集, 利用$TB(\theta,n)$的对称性知$U(-y_0,r)\subseteq\overline{TB(\theta,n)}$, 再由$TB(\theta,n)$是凸集知
\begin{align*}
U(\theta,r)\subseteq\frac{1}{2}U(y_0,r)+\frac{1}{2}U(-y_0,r)\subseteq\overline{TB(\theta,n)}.
\end{align*}
由$T$的齐次性, 取$\delta=\frac{r}{3n}$, 便有
\begin{align}\label{eq::FIOJWIOIOWIJOIJOWIOIOIIOIII}
U\left( \theta ,3\delta \right) =U\left( \theta ,\frac{r}{n} \right) =\frac{1}{n}U\left( \theta ,r \right) \subseteq \frac{1}{n}\overline{TB\left( \theta ,n \right) }=\overline{TB\left( \theta ,1 \right) }.
\end{align}
\begin{figure}[H]
\centering
\includegraphics[scale=0.3]{泛函分析图2.3.2.png}
\caption{}
\label{figure:2.3.2}
\end{figure}

\item 证明: $TB(\theta,1)\supseteq U(\theta,\delta)$. $\forall y_0\in U(\theta,\delta)$, 要证$\exists x_0\in B(\theta,1)$, 使得$Tx_0=y_0$, 即求方程$Tx=y_0$在$B(\theta,1)$内的一个解$x_0$, 我们用逐次逼近法.

对$y_0\in U(\theta,\delta)$,由\eqref{eq::FIOJWIOIOWIJOIJOWIOIOIIOIII}式知$3y_0\in U(\theta,3\delta)\subseteq \overline{TB(\theta,1)}$,再由$\overline{TB(\theta,1)}$是闭集知,$\exists z_1\in B(\theta,1)$,使得
\begin{align*}
\|3y_0-Tz_1\|<\delta.
\end{align*}
取$x_1=\frac{z_1}{3}$,则$\|y_0-Tx_1\|<\frac{\delta}{3}$;

对$y_1=y_0-Tx_1\in U\left(\theta,\frac{\delta}{3}\right)$, 同理可知$\exists x_2\in B\left(\theta,\frac{1}{3^2}\right)$, 使得$\|y_1-Tx_2\|<\frac{\delta}{3^2}$;

……………………

对$y_n=y_{n-1}-Tx_n\in U\left(\theta,\frac{\delta}{3^n}\right)$, 同理可知 $\exists x_{n+1}\in B\left(\theta,\frac{1}{3^{n+1}}\right)$, 使得$\|y_n-Tx_{n+1}\|<\frac{\delta}{3^{n+1}}$;

于是
\begin{align}\label{eqFJWFHOHIJW343}
\sum_{n=1}^{\infty}{\left\| x_n \right\|}<\sum_{n=1}^{\infty}{\frac{1}{3^n}}\leqslant \frac{1}{2},
\end{align}
令$S_n\triangleq\sum\limits_{i=1}^{n}x_i,\forall n\in\mathbb{N}$,则$S_n\in D(T),\forall n\in \mathbb{N}$.并且由\eqref{eqFJWFHOHIJW343}知,对$\forall n,p\in \mathbb{N}$
\begin{align*}
\left\| S_{n+p}-S_n \right\| =\left\| \sum_{k=n+1}^{n+p}{x_k} \right\| \leqslant \sum_{k=n+1}^{n+p}{\left\| x_k \right\|}\rightarrow 0\quad \left( n,p\rightarrow \infty \right) .
\end{align*}
故$\{S_n\}$是$\mathscr{X}$中的基本列,又$\mathscr{X}$是完备的,故可设$S_n\to x_0\triangleq\sum\limits_{i=1}^{\infty}x_i(n\to\infty)$,则由\eqref{eqFJWFHOHIJW343}知$x_0\in B(\theta,1)$.
而
\begin{align*}
\|y_n\|&=\|y_{n-1}-Tx_n\|=\cdots=\left\|y_0-T(x_1+x_2+\cdots+x_n)\right\|<\frac{\delta}{3^n}\quad(\forall n\in\mathbb{N}),
\end{align*}
即得$TS_n\to y_0\quad(n\to\infty).$
又因为$T$是连续的且$S_n\to x_0(n\to\infty)$, 所以再由极限的唯一性知$Tx_0=y_0,$
即得$U(\theta,\delta)\subseteq TB(\theta,1)$.
\end{enumerate}
\end{proof}

\begin{corollary}[开映射定理的推论]\label{corollary:开映射--有界原像}
设 $\mathscr{X},\mathscr{Y}$ 是Banach空间，$A\in\mathscr{L}(\mathscr{X},\mathscr{Y})$ 是满射。
则存在常数 $M>0$，使得对每个 $y\in\mathscr{Y}$，都存在 $x\in\mathscr{X}$ 满足
\[
Ax = y \quad\text{且}\quad \|x\| \leqslant M\|y\|.
\]
\end{corollary}
\begin{proof}
由\hyperref[theorem:泛函分析--开映射定理]{开映射定理}知$A$是开映射。
故存在 $\delta>0$ 使得
\[
B_{\mathscr{Y}}(0,\delta)\subseteq A\bigl(B_{\mathscr{X}}(0,1)\bigr).
\]
对任意 $y\in\mathscr{Y}\setminus\{0\}$，向量 $\frac{\delta}{2\|y\|}y$ 的范数为 $\frac{\delta}{2}<\delta$，由$A$满射知存在 $\tilde{x}\in B_{\mathscr{X}}(0,1)$ 满足
\[
A\tilde{x}=\frac{\delta}{2\|y\|}y.
\]
令 $x=\frac{2\|y\|}{\delta}\tilde{x}$，则 $Ax=y$ 且
\[
\|x\|=\frac{2\|y\|}{\delta}\|\tilde{x}\|<\frac{2}{\delta}\|y\|.
\]
另外,当$y=0$时,取$x=0$上式仍成立.取$M=\frac{2}{\delta}$，我们得到
\begin{align*}
\forall y\in\mathscr{Y},\ \exists x\in\mathscr{X}\text{ 满足 } Ax=y \text{ 且 } \|x\|\leqslant M\|y\|.
\end{align*}
\end{proof}

\begin{theorem}[Banach逆算子定理]\label{theorem:泛函分析--Banach逆算子定理}
设$\mathscr{X},\mathscr{Y}$是$B$空间. 若$T\in\mathscr{L}(\mathscr{X},\mathscr{Y})$, 它既是单射又是满射, 那么$T^{-1}\in\mathscr{L}(\mathscr{Y},\mathscr{X})$.
\end{theorem}
\begin{remark}
定理中的$B$空间$\mathscr{X},\mathscr{Y}$可以换成更一般的$F$空间, 但证明需稍做修改. 参看关肇直、张恭庆、冯德兴所著《线性泛函分析入门》(上海科学技术出版社, 1979)的第二章§2.
\end{remark}
\begin{proof}
依\hyperref[theorem:泛函分析--开映射定理]{开映射定理}证明中的第(3)部分, $\exists \delta>0$,使得$U(\theta,1)\subseteq TB\left(\theta,\frac{1}{\delta}\right)$, 即
\begin{align*}
T^{-1}U(\theta ,1)\subseteq B\left( \theta ,\frac{1}{\delta} \right) \text{或}\parallel T^{-1}y\parallel <\frac{1}{\delta}\quad (\forall y\in \mathscr{Y} ,\parallel y\parallel <1).
\end{align*}
特别地, 由范数的齐次性, $\forall y\in\mathscr{Y},\forall\varepsilon>0$, 有$\left\| \frac{y}{\left( 1+\varepsilon \right) \left\| y \right\|} \right\| <1$,进而由上式知
\begin{align*}
\left\| T^{-1}\frac{y}{\left( 1+\varepsilon \right) \left\| y \right\|} \right\| <\frac{1}{\delta}\Longrightarrow \left\| T^{-1}y \right\| <\frac{1+\varepsilon}{\delta}\left\| y \right\| .
\end{align*}
令$\varepsilon\to0$得
\begin{align*}
\|T^{-1}y\|\leqslant\frac{1}{\delta}\|y\|\quad(\forall y\in\mathscr{Y}).
\end{align*}
从而$T^{-1}\in\mathscr{L}(\mathscr{Y},\mathscr{X})$.
\end{proof}


\begin{definition}
设$\mathscr{X},\mathscr{Y}$是$B^*$空间,$T$是$\mathscr{X}\to\mathscr{Y}$的线性算子, $D(T)$是其定义域. 称$T$是\textbf{闭的}, 是指由$x_n\in D(T),x_n\to x$, 以及$Tx_n\to y$就能推出$x\in D(T)$, 而且$y=Tx$.
\end{definition}

\begin{proposition}\label{proposition:有界线性算子必是闭的}
\begin{enumerate}[(1)]
\item 设$\mathscr{X},\mathscr{Y}$是$B^*$空间,$T: \mathscr{X} \to  \mathscr{X}$ 为有界线性算子，则 $T$ 是闭算子。

\item 设 $\mathscr{X},\mathscr{Y}$ 是Banach空间，$T:D(T)\subseteq\mathscr{X}\to\mathscr{Y}$ 是闭线性算子且为单射，则其逆算子 $T^{-1}:R(T)\to D(T)$ 也是闭线性算子。
\end{enumerate}
\end{proposition}
\begin{proof}
\begin{enumerate}[(1)]
\item 设 $\{x_n\} \subset \mathscr{X}$ 满足 $x_n \to x$且 $Tx_n \to y$。由于 $T$ 有界，从而由\refpro{proposition:线性算子有界当且仅当连续}知$T$是连续的，故 $Tx_n \to Tx$。又极限唯一，所以 $y = Tx$。显然 $x \in \mathscr{X} = D(T)$。因此 $T$ 是闭的。

\item 任取序列 $\{y_n\}\subseteq R(T)$，满足
\begin{align*}
y_n\to y ,\quad T^{-1}y_n\to x.
\end{align*}
令 $x_n=T^{-1}y_n\in D(T)$，则 $Tx_n=y_n$. 于是有
\begin{align*}
x_n\to x,\quad Tx_n=y_n\to y.
\end{align*}
因 $T$ 是闭算子，由闭算子的序列定义即得 $x\in D(T)$ 且 $Tx=y$. 故 $y=Tx\in R(T)$ 且 $T^{-1}y=x$. 这就验证了 $T^{-1}$ 的闭性.
\end{enumerate}
\end{proof}

\begin{examples}
在$C[0,1]$上, $D(T)=C^1[0,1], T=\frac{\mathrm{d}}{\mathrm{d}t}$是一个闭线性算子.
\end{examples}
\begin{proof}
如果$x_n(t)\in C^1[0,1]$, 并且有
\begin{align*}
x_n\to x(C[0,1]),\quad \frac{\mathrm{d}x_n}{\mathrm{d}t}\to y(C[0,1]),
\end{align*}
则有
\begin{align*}
x_n(t)-x_n(0)\to\int_{0}^{t}y(\tau)\mathrm{d}\tau\quad(\forall t\in[0,1]),\\
x_n(t)-x_n(0)\to x(t)-x(0)\quad(\forall t\in[0,1]),
\end{align*}
即得
\begin{align*}
x(t)=x(0)+\int_{0}^{t}y(\tau)\mathrm{d}\tau\quad(\forall t\in[0,1]).
\end{align*}
因此, $x\in C^1[0,1]$, 且$\frac{\mathrm{d}x}{\mathrm{d}t}=y(t)$.
\end{proof}

\begin{theorem}
若$\mathscr{X},\mathscr{Y}$是$B$空间, $T$是$\mathscr{X}\to\mathscr{Y}$的一个闭线性算子, 满足$R(T)$是$\mathscr{Y}$中的第二纲集, 则$R(T)=\mathscr{Y}$并且$\forall\varepsilon>0,\exists\delta=\delta(\varepsilon)>0$, 使得$\forall y\in\mathscr{Y},\|y\|<\delta$必有$x\in D(T)$, 适合$\|x\|<\varepsilon$且$y=Tx$.
\end{theorem}
\begin{remark}
在这个定理中, $T\mathscr{X}$是第二纲集的假设是不可少的(满射及$\mathscr{Y}$的完备性保证了这一点). 因为有例子, 取$\mathscr{X}=\mathscr{Y}=C[0,1]$, 规定
\begin{align*}
(Tx)(t)=\int_{0}^{t}x(\tau)\mathrm{d}\tau\quad(\forall x\in\mathscr{X}).
\end{align*}
它显然是连续线性的, 但$T\mathscr{X}=\mathscr{Y}_0=\{y\in C^1[0,1]|y(0)=0\}$不是$C[0,1]$的第二纲集. 这时$T^{-1}=\frac{\mathrm{d}}{\mathrm{d}t}$在$C[0,1]$中不是连续的(即使以$C[0,1]$中的一个子集$\mathscr{Y}_0$作为$T^{-1}$的定义域, 也不连续). 事实上, $x_n(t)\triangleq\sin n\pi t$, 显然$\|x_n\|=1$, 但是
\begin{align*}
\left\|\frac{\mathrm{d}}{\mathrm{d}t}x_n(t)\right\|=n\pi\|\cos n\pi t\|=n\pi\to\infty\quad(\text{当}n\to\infty),
\end{align*}
其中$\|\cdot\|$表示$C[0,1]$空间中的范数. 然而, 若$\mathscr{Y}_0$按$C^1[0,1]$的范数$\|\cdot\|_1$则构成$B$空间, 这时$T^{-1}=\frac{\mathrm{d}}{\mathrm{d}t}$是有界的. 事实上,
\begin{align*}
\|T^{-1}y\|=\left\|\frac{\mathrm{d}}{\mathrm{d}t}y(t)\right\|\leqslant\|y\|_1\quad(\forall y\in\mathscr{Y}_0).
\end{align*}
\end{remark}
\begin{proof}
用$B(x_0,a),U(y_0,b)$分别表示$\mathscr{X},\mathscr{Y}$中的开球.令$B_n\triangleq B(\theta,n)\cap D(T)$,则
\begin{align*}
R(T)=\bigcup\limits_{n=1}^{\infty}TB_n,
\end{align*}
而$R(T)$是第二纲集,所以至少有一个$n\in\mathbb{N}$, 使得$TB_n$非疏, 即$\overline{TB_n}$至少含有一个内点. 因此$\exists U(y_0,r)\subseteq\overline{TB_n}$, 由$T$的线性性质知$TB_n$是一个对称凸集, 利用$TB_n$的对称性知$U(-y_0,r)\subseteq\overline{TB_n}$, 再由$TB_n$是凸集知
\begin{align*}
U(\theta,r)\subseteq\frac{1}{2}U(y_0,r)+\frac{1}{2}U(-y_0,r)\subseteq\overline{TB_n}.
\end{align*}
由$T$的齐次性, 取$\delta=\frac{r}{3n}$, 便有
\begin{align}\label{eq::FIOJWIOIOWIJOIJOWIOIOIIOIII1}
U\left( \theta ,3\delta \right) =U\left( \theta ,\frac{r}{n} \right) =\frac{1}{n}U\left( \theta ,r \right) \subseteq \frac{1}{n}\overline{TB_n }=\overline{TB_n }.
\end{align}
再利用逐次逼近法.

对$y_0\in U(\theta,\delta)$,由\eqref{eq::FIOJWIOIOWIJOIJOWIOIOIIOIII1}式知$3y_0\in U(\theta,3\delta)\subseteq \overline{TB_1}$,再由$\overline{TB_n}$是闭集知,$\exists z_1\in B_1$,使得
\begin{align*}
\|3y_0-Tz_1\|<\delta.
\end{align*}
取$x_1=\frac{z_1}{3}$,则$\|y_0-Tx_1\|<\frac{\delta}{3}$;

对$y_1=y_0-Tx_1\in U\left(\theta,\frac{\delta}{3}\right)$, 再由\eqref{eq::FIOJWIOIOWIJOIJOWIOIOIIOIII1}式同理可知$\exists x_2\in B_{\frac{1}{3^2}}$, 使得$\|y_1-Tx_2\|<\frac{\delta}{3^2}$;

……………………

对$y_n=y_{n-1}-Tx_n\in U\left(\theta,\frac{\delta}{3^n}\right)$, 再由\eqref{eq::FIOJWIOIOWIJOIJOWIOIOIIOIII1}式同理可知 $\exists x_{n+1}\in B_{\frac{1}{3^{n+1}}}$, 使得$\|y_n-Tx_{n+1}\|<\frac{\delta}{3^{n+1}}$;

于是
\begin{align}\label{eqFJWFHOHIJW3431}
\sum_{n=1}^{\infty}{\left\| x_n \right\|}<\sum_{n=1}^{\infty}{\frac{1}{3^n}}\leqslant \frac{1}{2},
\end{align}
令$S_n\triangleq\sum\limits_{i=1}^{n}x_i,\forall n\in\mathbb{N}$,则由\eqref{eqFJWFHOHIJW3431}知,对$\forall n,p\in \mathbb{N}$
\begin{align*}
\left\| S_{n+p}-S_n \right\| =\left\| \sum_{k=n+1}^{n+p}{x_k} \right\| \leqslant \sum_{k=n+1}^{n+p}{\left\| x_k \right\|}\rightarrow 0\quad \left( n,p\rightarrow \infty \right) .
\end{align*}
故$\{S_n\}$是$\mathscr{X}$中的基本列,又$\mathscr{X}$是完备的,故可设$S_n\to x_0\triangleq\sum\limits_{i=1}^{\infty}x_i(n\to\infty)$,则由\eqref{eqFJWFHOHIJW3431}知$x_0\in B_1$.
而
\begin{align*}
\|y_n\|&=\|y_{n-1}-Tx_n\|=\cdots=\left\|y_0-T(x_1+x_2+\cdots+x_n)\right\|<\frac{\delta}{3^n}\quad(\forall n\in\mathbb{N}),
\end{align*}
即得$TS_n\to y_0\quad(n\to\infty).$
又因为$S_n\in D(T)(\forall n\in \mathbb{N})$且$S_n\to x_0(n\to\infty)$, 所以由$T$是闭的知$Tx_0=y_0,$
即得
\begin{align}\label{eqHIUOFHJUW}
U(\theta,\delta)\subseteq TB_1=T(B(\theta,1)\cap D(T)).
\end{align}

先证$R(T)=\mathscr{Y}$.由\eqref{eqHIUOFHJUW}式知$\exists\delta>0$, 使得
\begin{align*}
U(\theta,\delta)\subseteq T(B(\theta,1)\cap D(T)).
\end{align*}
$\forall y\in\mathscr{Y}$, 不妨设$y\neq\theta$(显然$\theta\in R(T)$). $\forall0<\delta_1<\delta$, 按式,
\begin{align*}
\frac{\delta_1 y}{\|y\|}\in U(\theta,\delta)\implies\frac{\delta_1 y}{\|y\|}\in T(B(\theta,1)\cap D(T)).
\end{align*}
于是$\exists x\in B(\theta,1)\cap D(T)$, 使得
\begin{align*}
\frac{\delta_1 y}{\|y\|}=Tx\implies y=T\left(\frac{\|y\|}{\delta_1}x\right)\implies y\in R(T).
\end{align*}
故$\mathscr{Y}\subseteq R(T)$,又显然有$\mathscr{Y}\supseteq R(T)$,因此$\mathscr{Y} = R(T)$.

对$\forall \varepsilon>0$，令$\delta(\varepsilon)=\varepsilon\delta$，若$\|y\|<\delta(\varepsilon)=\varepsilon\delta$，则$\left\|\frac{y}{\varepsilon}\right\|<\delta$，即$\frac{y}{\varepsilon}\in U(\theta,\delta)$。
从而由\eqref{eqHIUOFHJUW}式知存在$x_0\in B(\theta,1)\cap D(T)$满足$Tx_0=\frac{y}{\varepsilon}$。取$x=\varepsilon x_0$，则$x\in D(T)$，$\|x\|<\varepsilon$，且$Tx=y$。
\end{proof}


\subsection{闭图像定理}

\begin{theorem}
设$T$是$B^*$空间$\mathscr{X}$到$B$空间$\mathscr{Y}$的连续线性算子, 那么$T$能唯一地延拓到$\overline{D(T)}$上成为连续线性算子$T_1$, 使得$T_1|_{D(T)}=T$, 且$\|T_1\|=\|T\|$.
\end{theorem}
\begin{proof}
任取$x\in\overline{D(T)},\exists x_n\in D(T),\lim\limits_{n\to\infty}x_n=x$, 依假设$T$在$D(T)$上连续, 由\refpro{proposition:线性算子有界当且仅当连续}知$T$有界, 即$\exists M>0$, 使得
\begin{align*}
\|Tx\|\leqslant M\|x\|\quad(\forall x\in D(T)).
\end{align*}
于是
\begin{align*}
\|Tx_{n+p}-Tx_n\|\leqslant M\|x_{n+p}-x_n\|.
\end{align*}
由此可见$\{Tx_n\}$是$\mathscr{Y}$中的基本列, 已设$\mathscr{Y}$完备, 所以$\exists y\in\mathscr{Y}$, 使得$Tx_n\to y=\lim_{n\to \infty}Tx_n$. 不难看出$y$仅依赖于$x$, 而与$D(T)$中$x_n$的选择无关. 因此, 可以定义$T_1:x\mapsto y= \lim_{n\to \infty}Tx_n$. 容易验证$T_1$是线性的, 还有$T_1|_{D(T)}=T$. 对$\forall x\in\overline{D(T)}$,存在$D(T)$中的序列$x_n\rightarrow x\left( n\rightarrow \infty \right)$.从而
\begin{align*}
\left\| T_1x \right\| =\underset{n\rightarrow \infty}{\lim}\left\| Tx_n \right\| \leqslant M\underset{n\rightarrow \infty}{\lim}\left\| x_n \right\| =M\left\| x \right\| .
\end{align*}
故$T_1$有界,由\refpro{proposition:线性算子有界当且仅当连续}知$T$连续.
\end{proof}

\begin{corollary}[等价范数定理]\label{corollary:泛函分析--等价范数定理}
设线性空间$\mathscr{X}$上有两个范数$\|\cdot\|_1$与$\|\cdot\|_2$. 如果$\mathscr{X}$关于这两个范数都构成$B$空间, 而且$\|\cdot\|_2$比$\|\cdot\|_1$强, 则$\|\cdot\|_2$与$\|\cdot\|_1$必等价.
\end{corollary}
\begin{proof}
考察恒同映射$I:\mathscr{X}\to\mathscr{X}$, 把它看成是由$(\mathscr{X},\|\cdot\|_2)\to(\mathscr{X},\|\cdot\|_1)$的线性算子, 由假设$\|\cdot\|_2$比$\|\cdot\|_1$强, 从而由\refpro{proposition:范数强的充要条件}知$\exists C>0$, 使得
\begin{align*}
\|Ix\|_1 \leqslant C\|x\|_2 \quad (\forall x\in\mathscr{X}).
\end{align*}
因此$I$是连续的, 由\refpro{proposition:线性算子有界当且仅当连续}知$I$有界.又注意到$I$既是单射又是满射,依\hyperref[theorem:泛函分析--Banach逆算子定理]{Banach逆算子定理}, $I$可逆且$I^{-1}$连续, 由\refpro{proposition:线性算子有界当且仅当连续}知$I^{-1}$有界,即有$M>0$, 使
\begin{align*}
\|I^{-1}x\|_2 \leqslant M\|x\|_1 \quad (\forall x\in\mathscr{X}).
\end{align*}
又因$Ix,I^{-1}x$与$x$是同一个元素, 所以
\begin{align*}
\frac{1}{C}\left\| x \right\| _1=\frac{1}{C}\left\| Ix \right\| _1\leqslant \left\| x \right\| _2=\left\| I^{-1}x \right\| _2\leqslant M\left\| x \right\| _1\quad (\forall x\in \mathscr{X} ).
\end{align*}
故由\refcor{corollary:范数等价的充要条件}知$\|\cdot\|_1$与$\|\cdot\|_2$等价.
\end{proof}

\begin{definition}
设$\mathscr{X},\mathscr{Y}$是线性空间,$T$是$\mathscr{X}\to\mathscr{Y}$的线性算子,则集合$G(T)\triangleq\{(x,Tx)\mid x\in D(T)\}$称为算子$T$的\textbf{图像}, 

设$\mathscr{X},\mathscr{Y}$是赋范线性空间,范数为$\|\cdot\|$,$T$是$\mathscr{X}\to\mathscr{Y}$的线性算子.引进另外一个范数$\|\cdot\|_G$如下:
\begin{align*}
\left\| x \right\| _G=\left\| x \right\| _{\mathscr{X}}+\left\| Tx \right\| _{\mathscr{Y}}\quad (\forall x\in D(T)).
\end{align*}
则$\|x\|_G$实际上是$(x,Tx)$在乘积空间$\mathscr{X}\times\mathscr{Y}$上的范数, 因此将$\|\cdot\|_G$称为\textbf{图模}. 
\end{definition}
\begin{remark}
也可将上述$\|x\|_G$视作$D(T)$空间上的范数,此时显然有$\|\cdot\|_G$比$\|\cdot\|_\mathscr{X}$和$\|\cdot\|_\mathscr{Y}$强.
\end{remark}

\begin{proposition}
设$\mathscr{X},\mathscr{Y}$是$B$空间,则算子$T$是闭的当且仅当$G(T)$按图模是闭的.
\end{proposition}
\begin{proof}
\end{proof}

\begin{theorem}[闭图像定理]\label{theorem:泛函分析--闭图像定理}
设$\mathscr{X},\mathscr{Y}$是$B$空间. 若$T$是$\mathscr{X}\to\mathscr{Y}$的闭线性算子, 并且$D(T)$是闭的, 则$T$是连续的.
\end{theorem}
\begin{proof}
因为$D(T)$是闭的, 所以由\rrefpro{proposition:完备度量空间的子空间等价于闭集}{proposition:完备度量空间的子空间等价于闭集-1}知$D(T)$作为$\mathscr{X}$的线性子空间可看成按$\mathscr{X}$中的范数$\|\cdot\|$构成$B$空间. 在$D(T)$上, 引进另外一个范数$\|\cdot\|_G$如下:
\begin{align*}
\|x\|_G = \|x\| + \|Tx\| \quad (\forall x\in D(T)).
\end{align*}
现在证明$(D(T),\|\cdot\|_G)$也是$B$空间, 事实上, 从
\begin{align*}
\|x_n - x_m\|_G = \|x_n - x_m\| + \|Tx_n - Tx_m\| \to 0 \quad (n,m\to\infty),
\end{align*}
可知$\exists x^*\in\mathscr{X}$与$y^*\in\mathscr{Y}$, 使得$x_n\to x^*$, 且$Tx_n\to y^*$. 根据$T$的闭性即得$y^* = Tx^*$, 从而$Tx_n\to Tx^*$. 因此$\|x_n - x^*\|_G\to 0(n\to \infty)$,故$(D(T),\|\cdot\|_G)$也是$B$空间.
又显然有$\|\cdot\|_G$比$\|\cdot\|$强, 根据\hyperref[corollary:泛函分析--等价范数定理]{等价范数定理}, $\|\cdot\|_G$与$\|\cdot\|$等价, 故$\exists M>0$, 使得
\begin{align*}
\|Tx\| \leqslant \|x\|_G \leqslant M\|x\| \quad (\forall x\in D(T)).
\end{align*}
故$T$有界,由\refpro{proposition:线性算子有界当且仅当连续}知$T$连续.
\end{proof}


\subsection{共鸣定理}

\begin{theorem}[共鸣定理或一致有界定理]\label{theorem:泛函分析--共鸣定理}
设$\mathscr{X}$是$B$空间, $\mathscr{Y}$是$B^*$空间, 如果$W\subseteq\mathscr{L}(\mathscr{X},\mathscr{Y})$, 使得
\begin{align*}
\sup_{A\in W}\|Ax\| < \infty \quad (\forall x\in\mathscr{X}),
\end{align*}
那么存在常数$M$, 使得$\|A\|\leqslant M(\forall A\in W)$.

逆否命题就是:如果$\sup_{A\in W}\|A\|=\infty$,则$\exists x_0\in\mathscr{X}$, 使得
\begin{align*}
\sup_{A\in W}\|Ax_0\| = \infty.
\end{align*}
\end{theorem}
\begin{remark}
条件: $\forall x\in\mathscr{X}$, $\sup_{A\in W}\|Ax\|<\infty$, 意味着$\forall x\in\mathscr{X}$, $\exists M_x>0$, 使得
\begin{align}
\|Ax\| \leqslant M_x\|x\| \quad (\forall A\in W).
\label{eq:HIFHIOW0-3823412}
\end{align}
而结论: $\|A\|\leqslant M(\forall A\in W)$, 则可看作是, 存在与$x$无关的常数$M$, 使得
\begin{align}
\|Ax\| \leqslant M\|x\| \quad (\forall A\in W).
\label{eq:HIFHIOW0-3823413}
\end{align}
式\eqref{eq:HIFHIOW0-3823412}意味着算子族$W$点点有界; 式\eqref{eq:HIFHIOW0-3823413}则意味着算子族$W$一致有界. 因此本定理给出条件保证点点有界蕴含一致有界, 故称“一致有界”定理. 另一方面, 如果我们从反面来叙述本定理将有: $\sup_{A\in W}\|A\|=\infty\implies\exists x_0\in\mathscr{X}$, 使得
\begin{align*}
\sup_{A\in W}\|Ax_0\| = \infty.
\end{align*}
因此本定理又有“共鸣定理”之称.
\end{remark}
\begin{proof}
$\forall x\in\mathscr{X}$, 定义
\begin{align*}
\|x\|_W = \|x\| + \sup_{A\in W}\|Ax\|.
\end{align*}
显然, $\|\cdot\|_W$是$\mathscr{X}$上的范数, 且强于$\|\cdot\|$. 下面证明$(\mathscr{X},\|\cdot\|_W)$完备. 事实上, 如果
\begin{align*}
\|x_m - x_n\| + \sup_{A\in W}\|A(x_m - x_n)\| \to 0 \quad (\text{当 }m,n\to\infty).
\end{align*}
由$\mathscr{X}$的完备性, $\exists x\in\mathscr{X}$, 使得$\|x_n - x\|\to 0$ (当$n\to\infty$), 又因为$\forall\varepsilon>0$, $\exists N=N(\varepsilon)$, 使得
\begin{align*}
\sup_{A\in W}\|Ax_m - Ax_n\| < \varepsilon \quad (\forall m,n\geqslant N).
\end{align*}
从而对$\forall A\in W$有$\|Ax_n - Ax\|\leqslant\varepsilon$ ($\forall n\geqslant N$). 于是
\begin{align*}
\|x_n - x\| + \sup_{A\in W}\|A(x_n - x)\| \to 0 \quad (\text{当 }n\to\infty),
\end{align*}
即$\|x_n - x\|_W\to 0$.故$(\mathscr{X},\|\cdot\|_W)$和$(\mathscr{X},\|\cdot\|)$都是$B$空间. 再根据\hyperref[corollary:泛函分析--等价范数定理]{等价范数定理}, $\|\cdot\|_W$与$\|\cdot\|$等价, 从而存在常数$M$, 使得
\begin{align*}
\sup\limits_{A\in W}\|Ax\| \leqslant M\|x\| \quad (\forall x\in \mathscr{X}) \Longrightarrow \sup\limits_{A\in W}\|A\| = \sup\limits_{\substack{A\in W\\ x\in \mathscr{X} \setminus \{ \theta \}}}\frac{\|Ax\|}{\|x\|} \leqslant M.
\end{align*}
由此立即推出$\|A\|\leqslant M(\forall A\in W)$.
\end{proof}

\begin{theorem}[Banach-Steinhaus定理]\label{theorem:泛函分析--Banach-Steinhaus定理}
设$\mathscr{X}$是$B$空间, $\mathscr{Y}$是$B^*$空间, $M$是$\mathscr{X}$的某个稠密子集. 若$A_n(n=1,2,\cdots),A\in\mathscr{L}(\mathscr{X},\mathscr{Y})$, 则$\forall x\in\mathscr{X}$都有
\begin{align*}
\lim_{n\to\infty}A_nx = Ax
\end{align*}
的充要条件是:
\begin{enumerate}[(1)]
\item $\|A_n\|(\forall n\in \mathbb{N})$有界;
\item $\lim_{n\to\infty}A_nx = Ax$对$\forall x\in M$成立.
\end{enumerate}
\end{theorem}
\begin{proof}

{\heiti 必要性:} 由假设可知结论(2)显然成立.记$W=\{A_n\}_{n=1}^{\infty}$,则由假设知
\begin{align*}
\mathop {\mathrm{sup}} \limits_{T\in W}\left\| Tx \right\| =\mathop {\mathrm{sup}} \limits_{n\in \mathbb{N}}\left\| A_nx \right\| \leqslant Ax+1\,\quad (\forall x\in \mathscr{X} ).
\end{align*}
再根据\hyperref[theorem:泛函分析--共鸣定理]{共鸣定理}知$\|A_n\|(\forall n\in \mathbb{N})$有界.

{\heiti 充分性:} 假定$\|A_n\|\leqslant C(\forall n\in\mathbb{N})$, 由$M$的稠密性知对$\forall x\in\mathscr{X}$及$\forall\varepsilon>0$, 取$y\in M$, 使得
\begin{align*}
\|x - y\| \leqslant \frac{\varepsilon}{4(\|A\| + C)},
\end{align*}
由\refpro{proposition:线性算子有界当且仅当连续}知$A_n$都连续,再结合上式知,存在$N'\in \mathbb{N}$,使得
\begin{align*}
\|A_nx - Ax\| &\leqslant \|A_nx - A_ny\| + \|A_ny - Ay\| + \|Ax - Ay\| \\
&< \frac{\varepsilon}{2} + \|A_ny - Ay\| \quad (\forall n>N').
\end{align*}
再取$N$足够大, 使得$\|A_ny - Ay\| < \frac{\varepsilon}{2}$ ($\forall n\geqslant N$), 便有
\begin{align*}
\|A_nx - Ax\| < \varepsilon \quad (\forall n\geqslant N).
\end{align*}
\end{proof}


\subsection{应用}

\begin{theorem}[Lax-Milgram定理]\label{theorem:泛函分析--Lax-Milgram定理}
设$a(x,y)$是Hilbert空间$\mathscr{X}$上的一个共轭双线性函数, 满足:
\begin{enumerate}[(1)]
\item $\exists M>0$, 使$|a(x,y)|\leqslant M\|x\|\cdot\|y\|$ ($\forall x,y\in\mathscr{X}$);
\item $\exists\delta>0$, 使$|a(x,x)|\geqslant\delta\|x\|^2$ ($\forall x\in\mathscr{X}$),
\end{enumerate}
那么必存在唯一的有连续逆的连续线性算子$A\in\mathscr{L}(\mathscr{X})$, 满足
\begin{align}
a(x,y) = (x,Ay) \quad (\forall x,y\in\mathscr{X}),
\label{eq:HIFHIOW0-3823417}
\end{align}
\begin{align}
\|A^{-1}\| \leqslant \frac{1}{\delta}.
\label{eq:HIFHIOW0-3823418}
\end{align}
\end{theorem}
\begin{proof}
依\refthe{theorem:泛函分析--定理2.2.2}, 适合\eqref{eq:HIFHIOW0-3823417}式的算子$A\in\mathscr{L}(\mathscr{X})$存在且唯一. 今证:
\begin{enumerate}[(i)]
\item $A$是单射. 若有$y_1,y_2\in\mathscr{X}$, 满足$Ay_1=Ay_2$, 则由\eqref{eq:HIFHIOW0-3823417}式可得
\begin{align*}
a(x,y_1) = a(x,y_2) \quad (\forall x\in\mathscr{X}),
\end{align*}
从而
\begin{align*}
a(x,y_1 - y_2) = 0 \quad (\forall x\in\mathscr{X}).
\end{align*}
特别取$x=y_1 - y_2$, 由条件(2)得
\begin{align*}
\delta\|(y_1-y_2)^2\|\leqslant |a(y_1 - y_2,y_1 - y_2)|=0\Longrightarrow y_1-y_2=\theta \Longrightarrow y_1=y_2.
\end{align*}

\item $A$是满射. 先证$R(A)$是闭的. 事实上, $\forall w\in\overline{R(A)}$, $\exists v_n\in\mathscr{X}$ ($n=1,2,\cdots$), 使得
\begin{align}
w = \lim_{n\to\infty}Av_n.
\label{eq:HIFHIOW0-3823419}
\end{align}
由条件(2)、\eqref{eq:HIFHIOW0-3823417}式以及\hyperref[theorem:泛函分析--Cauchy-Schwarz不等式]{Cauchy-Schwarz不等式}得
\begin{align*}
\delta\|v_{n+p} - v_n\|^2 &\leqslant |a(v_{n+p} - v_n, v_{n+p} - v_n)| = |(v_{n+p} - v_n, A(v_{n+p} - v_n))| \\
&\leqslant \|v_{n+p} - v_n\| \cdot \|Av_{n+p} - Av_n\| \quad (\forall n,p\in\mathbb{N}),
\end{align*}
即得
\begin{align*}
\|v_{n+p} - v_n\| \leqslant \frac{1}{\delta}\|Av_{n+p} - Av_n\| \to 0 \quad (\text{当 }n\to\infty, \forall p\in\mathbb{N}).
\end{align*}
从而$\{v_n\}$是基本列, 因此$\exists v^*\in\mathscr{X}$, 使得$v_n\to v^*$, 并由$A$的连续性和式\eqref{eq:HIFHIOW0-3823419}得$w=Av^*$, 即$w\in R(A)$. 于是$R(A)$闭.

再证$R(A)^\perp=\{\theta\}$. 倘若$w\in R(A)^\perp$, 则
\begin{align*}
(w,Av) = 0 \quad (\forall v\in\mathscr{X}),
\end{align*}
由\eqref{eq:HIFHIOW0-3823417}式知$a(w,v)=0$ ($\forall v\in\mathscr{X}$). 特别取$v=w$, 再利用条件(2)有
\begin{align*}
\delta\|w\|^2 \leqslant |a(w,w)| = 0,
\end{align*}
即得$w=\theta$. 再由\rrefpro{proposition:Hilbert空间稠密线性空间正交补为0}{proposition:Hilbert空间稠密线性空间正交补为0-2}知$\overline{R(A)}=H$,故$R(A)^\perp=\{\theta\}$. 又因$R(A)$为闭子空间，故$R(A)=H$. 所以$A$是满射.

\item 再利用\hyperref[theorem:泛函分析--Banach逆算子定理]{Banach逆算子定理}可知$A^{-1}\in\mathscr{L}(\mathscr{X})$. 又由条件(2)、\eqref{eq:HIFHIOW0-3823417}式以及\hyperref[theorem:泛函分析--Cauchy-Schwarz不等式]{Cauchy-Schwarz不等式}可得
\begin{align*}
\delta\|x\|^2 \leqslant |a(x,x)| = |(x,Ax)| \leqslant \|x\| \cdot \|Ax\|,
\end{align*}
所以$\delta\|x\| \leqslant \|Ax\|$ ($\forall x\in\mathscr{X}$), 即得\eqref{eq:HIFHIOW0-3823418}式.
\end{enumerate}
\end{proof}

\begin{definition}[相容性]
设 $\mathscr{X} $ 和 $\mathscr{Y} $ 为 $B$空间, $T\in\mathscr{L} (\mathscr{X} ,\mathscr{Y} )$，$\{T_n\}_{n\in\mathbb{N}}\subset\mathscr{L} (\mathscr{X} ,\mathscr{Y} )$。称 $\{T_n\}$ 是 $T$ 的\textbf{近似}（或$T$的近似格式$\{T_n\}$具有\textbf{相容性}），若对每个 $x\in\mathscr{X} $，
\[
\|Tx - T_nx\| \to 0 \quad (n\to\infty).
\]
\end{definition}

\begin{definition}[稳定性]
设 $\mathscr{X} $ 和 $\mathscr{Y} $ 为 $B$空间,  $T_n\in\mathscr{L} (\mathscr{X} ,\mathscr{Y} )$ 可逆,即存在逆算子 $T_n^{-1}\in\mathscr{L} (\mathscr{Y} ,\mathscr{X} )$.称近似格式 $\{T_n\}$ 具有\textbf{稳定性}，若存在常数 $C>0$ 使得
\[
\|T_n^{-1}\| \leqslant C \quad (\forall n\in\mathbb{N}).
\]
\end{definition}

\begin{theorem}[Lax定理]\label{theorem:泛函分析--Lax定理}
设 $\mathscr{X} $ 和 $\mathscr{Y} $ 为 $B$空间,  $T\in\mathscr{L} (\mathscr{X} ,\mathscr{Y} )$ 可逆，$\{T_n\}\subset\mathscr{L} (\mathscr{X} ,\mathscr{Y} )$ 可逆。假设$T$的近似格式$\{T_n\}$具有相容性,即对 $\forall x\in\mathscr{X} $，有
\begin{align}\label{eq::23ur8902IOJAI1}
T_nx\to Tx(n\to \infty).
\end{align}
则对 $\forall y\in\mathscr{Y} $，近似方程 $T_nx_n=y$ 的解 $x_n = T_n^{-1}y$ 收敛于原方程 $Tx=y$ 的解 $x = T^{-1}y$,即 $x_n\to x$ 于 $\mathscr{X} $当且仅当$\{T_n\}$ 具有稳定性,即存在常数 $C>0$ 使得
\begin{align}
\|T_n^{-1}\| \leqslant C \quad (\forall n\in\mathbb{N}).\label{eq::23ur8902IOJAI}
\end{align}
\end{theorem}
\begin{proof}

{\heiti 充分性:} 由式\eqref{eq::23ur8902IOJAI1}和式\eqref{eq::23ur8902IOJAI}, 我们得
\begin{align*}
\|x_n - x\| = \|T_n^{-1}y - T_n^{-1}T_nx\| 
\leqslant \|T_n^{-1}\| \cdot \|Tx - T_nx\| 
\leqslant C\|Tx - T_nx\| \to 0 \quad (n\to\infty).
\end{align*}

{\heiti 必要性:} $\forall y\in\mathscr{Y}$, 令$x_n=T_n^{-1}y$, $x=T^{-1}y$, 便有$x_n\to x$ ($n\to\infty$). 因此,
\begin{align*}
T_n^{-1}y \to T^{-1}y \quad (n\to\infty, \forall y\in\mathscr{Y}).
\end{align*}
由\hyperref[theorem:泛函分析--共鸣定理]{共鸣定理}, 立得$\|T_n^{-1}\|$有界.
\end{proof}

\begin{example}
设$\mathscr{X}$是$B$空间，$\mathscr{X}_0$是$\mathscr{X}$的闭子空间。映射$\varphi:\mathscr{X}\to\mathscr{X}/\mathscr{X}_0$定义为
\begin{align*}
\varphi:x\mapsto [x]\quad (\forall x\in\mathscr{X}),
\end{align*}
其中$[x]\in \mathscr{X}/\mathscr{X}_0$表示含$x$的商类。求证$\varphi$是开映射。
\end{example}
\begin{proof}

由于 $\mathscr{X}_0$ 是闭子空间，商空间 $\mathscr{X}/\mathscr{X}_0$ 上赋予商范数
\[
\|[x]\| = \inf_{z\in\mathscr{X}_0} \|x+z\| \quad (\forall\,x\in\mathscr{X}),
\]
由\cref{theorem:闭线性子空间的商空间完备}知商空间在这个范数下也是Banach空间。现证 $\varphi$ 是开映射。

设 $V\subset\mathscr{X}$ 为任意开集，下证 $\varphi(V)$ 是 $\mathscr{X}/\mathscr{X}_0$ 中开集。
任取 $[x]\in\varphi(V)$，则存在 $x\in V$。因为 $V$ 开，存在 $\delta>0$ 使得开球 $B(x,\delta)\subset V$。
今证 $B([x],\delta)\subset\varphi(V)$。对任意 $[x']\in B([x],\delta)$，即 $\|[x']-[x]\|<\delta$，由商范数定义有
\[
\inf_{z\in\mathscr{X}_0} \|x'-x+z\| < \delta.
\]
故存在 $z_0\in\mathscr{X}_0$ 使得 $\|(x'+z_0)-x\|<\delta$。从而 $x'+z_0\in B(x,\delta)\subset V$。
在商映射下，$[x']=[x'+z_0]\in\varphi(V)$。因此 $B([x],\delta)\subset\varphi(V)$，说明 $\varphi(V)$ 是开集。
由 $V$ 的任意性，$\varphi$ 是开映射.
\end{proof}

\begin{example}
设$\mathscr{X},\mathscr{Y}$是$B$空间，又设方程$Ux=y$对$\forall y\in\mathscr{Y}$有解$x\in\mathscr{X}$，其中$U\in\mathscr{L}(\mathscr{X},\mathscr{Y})$，并且$\exists m>0$，使得
\begin{align*}
\|Ux\|\geqslant m\|x\|\quad (\forall x\in\mathscr{X}).
\end{align*}
求证：$U$有连续逆$U^{-1}$，并且$\|U^{-1}\|\leqslant \frac{1}{m}$。
\end{example}
\begin{proof}
由条件 $\|Ux\|\geqslant m\|x\|$ 知，若 $Ux=0$，则 $\|x\|=0$，故 $x=0$，所以 $U$ 是单射。
又由方程 $Ux=y$ 对任意 $y\in\mathscr{Y}$ 有解知 $U$ 是满射。
从而 $U$ 是双射，存在逆映射 $U^{-1}:\mathscr{Y}\to\mathscr{X}$，且易见 $U^{-1}$ 是线性的。
对任意 $y\in\mathscr{Y}$，令 $x=U^{-1}y$，则 $y=Ux$，由所给不等式得
\[
\|y\| = \|Ux\| \geqslant m\|x\| = m\|U^{-1}y\|,
\]
因此
\[
\|U^{-1}y\| \leqslant \frac{1}{m}\|y\| \quad (\forall y\in\mathscr{Y}).
\]
这表明 $U^{-1}$ 是有界线性算子，故由\cref{proposition:线性算子有界当且仅当连续}知$U^{-1}$ 连续.并且 $\|U^{-1}\|\leqslant \frac{1}{m}$.
\end{proof}

\begin{example}
设$\mathscr{X},\mathscr{Y}$是$B^*$空间，$D$是$\mathscr{X}$的线性子空间，并且$A:D\to\mathscr{Y}$是线性映射。求证：
\begin{enumerate}[(1)]
\item 如果$A$连续且$D$是闭的，那么$A$是闭算子；
\item 如果$A$连续且是闭算子，那么$\mathscr{Y}$完备蕴含$D$闭；
\item 如果$A$是单射的闭算子，那么$A^{-1}$也是闭算子；
\item 如果$\mathscr{X}$完备，$A$是单射的闭算子，$R(A)$在$\mathscr{Y}$中稠密，并且$A^{-1}$连续，那么$R(A)=\mathscr{Y}$。
\end{enumerate}
\end{example}
\begin{proof}
\begin{enumerate}[(1)]
\item 设 $\{x_n\}\subset D$ 满足 $x_n\to x\in\mathscr{X}$ 且 $Ax_n\to y\in\mathscr{Y}$。
由于 $D$ 是闭的，有 $x\in D$；又 $A$ 连续，故 $Ax_n\to Ax$。
极限唯一性给出 $y=Ax$，因此 $A$ 是闭算子。

\item 设 $\{x_n\}\subset D$ 且 $x_n\to x\in\mathscr{X}$,则$\{x_n\}$是$\mathscr{X}$中的Cauchy列.因 $A$ 连续，所以由\cref{proposition:线性算子有界当且仅当连续}知$A$有界,进而存在 $c>0$ 使
$\|Ax\|\leqslant c\|x\|$ 对所有 $x\in D$ 成立。于是
\begin{align*}
\|Ax_n-Ax_m\|\leqslant c\|x_n-x_m\|.
\end{align*}
又$\mathscr{X}$中的Cauchy列,故$\{Ax_n\}$ 为 $\mathscr{Y}$ 中 Cauchy 列。
由$\mathscr{Y}$ 完备，故存在 $y\in\mathscr{Y}$ 使 $Ax_n\to y$。
由 $A$ 是闭算子，$x_n\to x,\;Ax_n\to y$ 推出 $x\in D$，所以 $D$ 闭。

\item 因为$A$是单射,且$A$是$D$到$R(A)$的满射,所以存在逆映射$A^{-1}:R(A)\to D$.设 $\{y_n\}\subset R(A)$ 满足 $y_n\to y$ 且 $A^{-1}y_n\to x$。
令 $x_n=A^{-1}y_n\in D$，则 $Ax_n=y_n$。已知 $x_n\to x,\;Ax_n=y_n\to y$，
$A$ 是闭算子推出 $x\in D$ 且 $Ax=y$。故 $y\in R(A)$ 且 $A^{-1}y=x$，即 $A^{-1}$ 是闭算子。

\item 因为$A$是单射,且$A$是$D$到$R(A)$的满射,所以存在逆映射$A^{-1}:R(A)\to D$.只需证 $R(A)$ 闭。设 $\{y_n\}\subset R(A)$ 且 $y_n\to y\in\mathscr{Y}$,则$\{y_n\}$是$\mathscr{Y}$中的Cauchy列.
令 $x_n=A^{-1}y_n\in D$。因$A^{-1}$ 连续,由\cref{proposition:线性算子有界当且仅当连续}知$A^{-1}$有界,进而存在 $c>0$ 使得
\begin{align*}
\|x_n-x_m\|\leqslant c\|y_n-y_m\|.
\end{align*}
又$\{y_n\}$是$\mathscr{Y}$中的Cauchy列,故 $\{x_n\}$ 是 $\mathscr{X}$ 中 Cauchy 列。
由$\mathscr{X}$ 完备知,存在 $x\in\mathscr{X}$ 使 $x_n\to x$。
此时 $x_n\to x,\;Ax_n=y_n\to y$，$A$ 是闭算子推出 $x\in D$ 且 $Ax=y$。
因此 $y\in R(A)$，$R(A)$ 闭。又已知 $R(A)$ 在 $\mathscr{Y}$ 中稠密，故 $R(A)=\overline{R(A)}=\mathscr{Y}$。
\end{enumerate}
\end{proof}

\begin{example}
用\hyperref[corollary:泛函分析--等价范数定理]{等价范数定理}证明：$(C[0,1],\|\cdot\|_1)$不是$B$空间，其中$\|f\|_1=\int_{0}^{1}|f(t)|\mathrm{d}t,\forall f\in C[0,1]$。
\end{example}
\begin{proof}

用反证法。由\subcref{proposition:常见F空间}{proposition:常见F空间-1}知$(C[0,1],\|\cdot\|_\infty)$ 是 $B$ 空间，且对任意 $f\in C[0,1]$ 有
\[
\|f\|_1 = \int_0^1 |f(t)|\,\mathrm{d}t \leqslant \max_{t\in[0,1]}|f(t)| = \|f\|_\infty,
\]
即 $\|\cdot\|_\infty$ 强于 $\|\cdot\|_1$。
假设 $(C[0,1],\|\cdot\|_1)$ 也是 $B$ 空间。由\hyperref[corollary:泛函分析--等价范数定理]{等价范数定理}，
$\|\cdot\|_1$和$\|\cdot\|_{\infty}$等价.
从而存在常数 $c>0$，使得
\[
\|f\|_\infty \leqslant c\|f\|_1 \quad (\forall f\in C[0,1]).
\]

下证此式不可能成立。对每个正整数 $n\geqslant 2$，构造 $f_n\in C[0,1]$ 如下：
\begin{align*}
f_n(t) =
\begin{cases}
n^2t, & 0 \leqslant t \leqslant \frac{1}{n}, \\[2pt]
2n - n^2t, & \frac{1}{n} < t \leqslant \frac{2}{n}, \\[2pt]
0, & \frac{2}{n} < t \leqslant 1.
\end{cases}
\end{align*}
则 $\|f_n\|_\infty = n$，而
\[
\|f_n\|_1 = \int_0^{\frac{2}{n}} f_n(t)\,\mathrm{d}t = 1.
\]
于是 $\frac{\|f_n\|_\infty}{\|f_n\|_1} = n \to \infty\ (n\to\infty)$，
这与存在 $c>0$ 使得 $\|f\|_\infty \leqslant c\|f\|_1$ 对所有连续函数成立相矛盾!
故 $(C[0,1],\|\cdot\|_1)$ 不是 $B$ 空间.
\end{proof}

\begin{lemma}[Gelfand引理]\label{lemma:泛函分析--Gelfand引理}
设$\mathscr{X}$是$B$空间，$p:\mathscr{X}\to\mathbb{R}$满足
\begin{enumerate}[(1)]
\item $p(x)\geqslant 0\quad (\forall x\in\mathscr{X})$;
\item $p(\lambda x)=\lambda p(x)\quad (\forall\lambda>0,\forall x\in\mathscr{X})$;
\item $p(x_1+x_2)\leqslant p(x_1)+p(x_2)\quad (\forall x_1,x_2\in\mathscr{X})$;
\item 当$x_n\to x$时，$\varliminf\limits_{n\to\infty}p(x_n)\geqslant p(x)$.
\end{enumerate}
求证：$\exists M>0$，使得$p(x)\leqslant M\|x\|,\forall x\in\mathscr{X}$。
\end{lemma}
\begin{proof}

对每个 $n\in\mathbb{N}$，令
\[
F_n = \{x\in\mathscr{X} : p(x) \leqslant n\}.
\]
首先证明 $F_n$ 是闭集。设 $\{x_k\}\subset F_n$ 且 $x_k\to x$，则 $p(x_k)\leqslant n$。
由条件 (4) 知
\[
p(x) \leqslant \varliminf_{k\to\infty} p(x_k) \leqslant n,
\]
故 $x\in F_n$，因此 $F_n$ 为闭集。

由条件 (1) 知 $p(x)\geqslant 0$，且由$p$的陪域是$\mathbb{R}$知$p(x)<\infty$，故每个 $x$ 必属于某个 $F_n$，
从而 $\mathscr{X} = \bigcup_{n=1}^\infty F_n$。
因 $\mathscr{X}$ 是 Banach 空间，由\hyperref[theorem:泛函分析--Baire纲定理]{Baire纲定理}，存在 $N$ 使得 $F_N$ 有内点，
即存在 $x_0\in\mathscr{X}$ 和 $r>0$ 使得
\begin{align}\label{eq:ball}
B(x_0, r) \subseteq F_N.
\end{align}

对 $-x_0$，由$-x_0\in \mathscr{X} = \bigcup_{n=1}^\infty F_n$知,存在 $M$ 使得 $-x_0\in F_M$，即 $p(-x_0)\leqslant M$。
现任取 $y\in B(0,r)$。由 \eqref{eq:ball}式有 $y+x_0\in B(x_0,r)\subseteq F_N$，
故 $p(y+x_0)\leqslant N$。利用条件 (3) 的次可加性，
\[
p(y) \leqslant p(y+x_0) + p(-x_0) \leqslant N + M.
\]
因此 $B(0,r) \subseteq F_{N+M}$，即当 $\|x\|<r$ 时 $p(x)\leqslant N+M$。

对任意 $x\in\mathscr{X}\setminus\{0\}$，令 $t = \frac{r}{2\|x\|}>0$，
则 $\|tx\| = \frac{r}{2} < r$，于是 $p(tx) \leqslant N+M$。
由条件 (2) 的正齐次性知$p(tx)=t\,p(x)$，从而
\[
p(x) =\frac{p(tx)}{t}\leqslant \frac{N+M}{t} = \frac{2(N+M)}{r}\|x\|.
\]
此式对 $x=0$ 显然成立。
取 $M_0 = \frac{2(N+M)}{r}$ 即得
\[
p(x) \leqslant M_0 \|x\| \quad (\forall x\in\mathscr{X}),
\]
结论得证.
\end{proof}

\begin{example}
用 \hyperref[lemma:泛函分析--Gelfand引理]{Gelfand引理} 证明 \hyperref[theorem:泛函分析--共鸣定理]{共鸣定理}.
\end{example}
\begin{proof}
设 $\mathscr{X}$ 是 $B$ 空间, $\mathscr{Y}$ 是 $B^{*}$ 空间, 算子族 $\{A_{\alpha}\}_{\alpha\in\Lambda}\subseteq\mathscr{L}(\mathscr{X},\mathscr{Y})$ 满足
\begin{align*}
\sup\limits_{\alpha\in\Lambda}\|A_{\alpha}x\|<\infty\quad (\forall x\in\mathscr{X}).
\end{align*}
定义函数 $p:\mathscr{X}\to\mathbb{R}$ 为 $p(x)=\sup\limits_{\alpha\in\Lambda}\|A_{\alpha}x\|$. 下证 $p$ 满足 \hyperref[lemma:泛函分析--Gelfand引理]{Gelfand引理}的四个条件.
\begin{enumerate}[(1)]
\item 显然 $p(x)\geqslant 0$ 对任意 $x\in\mathscr{X}$ 成立.
\item 对任意 $\lambda>0$ 和 $x\in\mathscr{X}$,
\begin{align*}
p(\lambda x)=\sup\limits_{\alpha\in\Lambda}\|A_{\alpha}(\lambda x)\|
=\lambda\sup\limits_{\alpha\in\Lambda}\|A_{\alpha}x\|=\lambda p(x).
\end{align*}
\item 对任意 $x_{1},x_{2}\in\mathscr{X}$,
\begin{align*}
p(x_{1}+x_{2})
&=\sup\limits_{\alpha\in\Lambda}\|A_{\alpha}(x_{1}+x_{2})\|
\leqslant\sup\limits_{\alpha\in\Lambda}\bigl(\|A_{\alpha}x_{1}\|+\|A_{\alpha}x_{2}\|\bigr)\\
&\leqslant\sup\limits_{\alpha\in\Lambda}\|A_{\alpha}x_{1}\|+\sup\limits_{\alpha\in\Lambda}\|A_{\alpha}x_{2}\|
=p(x_{1})+p(x_{2}).
\end{align*}

\item 设 $\{x_{n}\}\subset\mathscr{X}$ 且 $x_{n}\to x\in\mathscr{X}$. 对每个固定的 $\alpha\in\Lambda$, 由 $A_{\alpha}$ 的连续性得 $\|A_{\alpha}x_{n}\|\to\|A_{\alpha}x\|$. 于是
\begin{align*}
\mathop {\mathrm{sup}} \limits_{\alpha \in \Lambda}\parallel A_{\alpha}x_n\parallel \geqslant \parallel A_{\alpha}x_n\parallel \Longrightarrow \varliminf\limits_{n\to\infty}p(x_{n})
&=\varliminf\limits_{n\to\infty}\sup\limits_{\alpha\in\Lambda}\|A_{\alpha}x_{n}\|
\geqslant\sup\limits_{\alpha\in\Lambda}\varliminf\limits_{n\to\infty}\|A_{\alpha}x_{n}\|
=\sup\limits_{\alpha\in\Lambda}\|A_{\alpha}x\|=p(x).
\end{align*}
\end{enumerate}
故由\hyperref[lemma:泛函分析--Gelfand引理]{Gelfand引理}知, 存在常数 $M>0$ 使得
\begin{align*}
p(x)\leqslant M\|x\|\quad (\forall x\in\mathscr{X}).
\end{align*}
由此即得 $\|A_{\alpha}x\|\leqslant M\|x\|$ 对一切 $\alpha\in\Lambda$ 及 $x\in\mathscr{X}$ 成立, 从而 $\sup\limits_{\alpha\in\Lambda}\|A_{\alpha}\|\leqslant M<\infty$.
\end{proof}

\begin{example}
设 $\mathscr{X}$ 和 $\mathscr{Y}$ 是 $B$ 空间, $A_n\in\mathscr{L}(\mathscr{X},\mathscr{Y})\ (n=1,2,\cdots)$, 又对 $\forall x\in\mathscr{X}$, $\{A_nx\}$ 在 $\mathscr{Y}$ 中收敛. 求证: $\exists A\in\mathscr{L}(\mathscr{X},\mathscr{Y})$, 使得
\begin{align*}
A_nx\to Ax\quad (\forall x\in\mathscr{X}),
\end{align*}
并且 $\|A\|\leqslant \varliminf\limits_{n\to\infty}\|A_n\|$.
\end{example}
\begin{proof}
定义算子 $A:\mathscr{X}\to\mathscr{Y}$ 为
\begin{align*}
Ax=\lim\limits_{n\to\infty}A_nx\quad (\forall x\in\mathscr{X}).
\end{align*}
易见 $A$ 是线性算子. 对每个 $x\in\mathscr{X}$, 序列 $\{A_nx\}$ 收敛, 故有界. 由\hyperref[theorem:泛函分析--共鸣定理]{共鸣定理}, 存在常数 $M>0$ 使得 
\begin{align*}
|A_n\|\leqslant M,\quad \forall n\in \mathbb{N}.
\end{align*}
对任意 $x\in\mathscr{X}$,由\cref{proposition:有界线性算子的基本不等式}得
\begin{align*}
\left\| A_nx \right\| \leqslant \left\| A_n \right\| \left\| x \right\| \Longrightarrow \|Ax\| = \lim\limits_{n\to\infty}\|A_nx\|
\leqslant \varliminf\limits_{n\to\infty}\bigl(\|A_n\|\|x\|\bigr)
= \Bigl(\varliminf\limits_{n\to\infty}\|A_n\|\Bigr)\|x\|.
\end{align*}
因此 $A$ 是有界算子, 且 $\|A\|\leqslant \varliminf\limits_{n\to\infty}\|A_n\|$. 特别地, 若极限 $\lim\limits_{n\to\infty}\|A_n\|$ 存在, 则 $\|A\|\leqslant \lim\limits_{n\to\infty}\|A_n\|$.
\end{proof}

\begin{example}
设 $1<p<\infty$, 并且 $\frac{1}{p}+\frac{1}{q}=1$. 如果序列 $\{\alpha_k\}$ 使得对 $\forall x=\{\xi_k\}\in l^p$ 保证 $\sum\limits_{k=1}^{\infty}\alpha_k\xi_k$ 收敛, 求证:
\begin{enumerate}[(1)]
\item $\{\alpha_k\}\in l^q$.

\item 若 $f:x\mapsto\sum\limits_{k=1}^{\infty}\alpha_k\xi_k$, 则 $f$ 作为 $l^p$ 上的线性泛函, 有
\begin{align*}
\|f\|=\left(\sum\limits_{k=1}^{\infty}|\alpha_k|^q\right)^{\frac{1}{q}}.
\end{align*}
\end{enumerate}
\end{example}
\begin{proof}
\begin{enumerate}[(1)]
\item 对每个 $n\in\mathbb{N}$, 定义线性泛函 $f_n:l^p\to\mathbb{R}$ (或 $\mathbb{C}$) 为
\begin{align*}
f_n(x)=\sum\limits_{k=1}^{n}\alpha_k\xi_k,\quad x=\{\xi_k\}\in l^p.
\end{align*}
由\hyperref[Basis of Analytics-theorem:离散的Hölder不等式]{H\"{o}lder不等式}得
\begin{align}\label{eq::jsijjjj-1}
|f_n(x)|\leqslant \sum\limits_{k=1}^{n}|\alpha_k\xi_k|\leqslant \Bigl(\sum\limits_{k=1}^{n}|\alpha_k|^q\Bigr)^{\frac{1}{q}}\|x\|_p\Longrightarrow \|f_n\| = \underset{\left\| x \right\| _p=1}{\mathrm{sup}}\left| f_n\left( x \right) \right|\leqslant \Bigl( \sum_{k=1}^n{|\alpha _k|^q} \Bigr) ^{\frac{1}{q}}. 
\end{align}
取$x_0=\{\xi_k^0\},\xi_k^0=\varepsilon_k^0|\alpha_k|^{q-1}$,其中$\varepsilon _k^0=\begin{cases}
\frac{\overline{\alpha _k}}{|\alpha _k|},&\alpha _k\ne 0,\\
0,&\alpha _k=0\\
\end{cases}$.则此时有
\begin{align*}
\frac{\left| f_n\left( x_0 \right) \right|}{\left\| x_0 \right\| _p}=\frac{\sum\limits_{k=1}^n{\alpha _k\varepsilon _k^0\left| \alpha _k \right|^{q-1}}}{\left( \sum\limits_{k=1}^n{\left| \varepsilon _k^0\left| \alpha _k \right|^{q-1} \right|^p} \right) ^{\frac{1}{p}}}=\frac{\sum\limits_{k=1}^n{\left| \alpha _k \right|^q}}{\left( \sum\limits_{k=1}^n{\left| \alpha _k \right|^{\left( q-1 \right) p}} \right) ^{\frac{1}{p}}}=\frac{\sum\limits_{k=1}^n{\left| \alpha _k \right|^q}}{\left( \sum\limits_{k=1}^n{\left| \alpha _k \right|^q} \right) ^{\frac{1}{p}}}=\left( \sum\limits_{k=1}^n{\left| \alpha _k \right|^q} \right) ^{\frac{1}{q}}.
\end{align*}
进而
\begin{align}\label{eq::jsijjjj-2}
\left\| f_n \right\| =\underset{x\in l^p}{\mathrm{sup}}\frac{\left| f_n\left( x \right) \right|}{\left\| x \right\| _p}\geqslant \left( \sum\limits_{k=1}^n{\left| \alpha _k \right|^q} \right) ^{\frac{1}{q}}.
\end{align}
因此由\eqref{eq::jsijjjj-1}\eqref{eq::jsijjjj-2}式知
\begin{align}\label{eq:fnorm}
\|f_n\|=\Bigl(\sum\limits_{k=1}^{n}|\alpha_k|^q\Bigr)^{\frac{1}{q}}.
\end{align}
依题设, 对每个 $x\in l^p$, $\lim\limits_{n\to\infty}f_n(x)=\sum\limits_{k=1}^{\infty}\alpha_k\xi_k$ 存在, 从而 $\{|f_n(x)|\}$ 有界. 由\hyperref[theorem:泛函分析--共鸣定理]{共鸣定理}, 存在常数 $M>0$ 使得 $\|f_n\|\leqslant M$ 对一切 $n$ 成立. 根据 \eqref{eq:fnorm}式即得
\begin{align*}
\Bigl(\sum\limits_{k=1}^{n}|\alpha_k|^q\Bigr)^{\frac{1}{q}}\leqslant M\quad \forall n\in \mathbb{N}.
\end{align*}
令 $n\to\infty$ 便知 $\{\alpha_k\}\in l^q$, 且 $\|\alpha\|_q\triangleq\bigl(\sum_{k=1}^{\infty}|\alpha_k|^q\bigr)^{\frac{1}{q}}\leqslant M$.

\item 现在定义 $f(x)=\sum\limits_{k=1}^{\infty}\alpha_k\xi_k$. 由 \hyperref[Real Analysis-theorem:实变函数--Hölder不等式]{H\"{o}lder不等式}得
\begin{align*}
|f(x)|\leqslant \sum\limits_{k=1}^{\infty}|\alpha_k\xi_k|
\leqslant \|\alpha\|_q\|x\|_p\quad (\forall x\in l^p),
\end{align*}
因此$f$0有界,所以$f\in (l^p)^*$ 且 $\|f\|\leqslant \|\alpha\|_q$.

再证反向不等式. 对每个 $n$, 构造 $x^{(n)}=\{\xi_k^{(n)}\}_{k=1}^{\infty}$ 如下:
\begin{align*}
\xi_k^{(n)}=
\begin{cases}
\frac{\varepsilon_k|\alpha_k|^{q-1}}{\Bigl(\sum\limits_{j=1}^{n}|\alpha_j|^q\Bigr)^{\frac{1}{p}}}, & 1\leqslant k\leqslant n,\\[4pt]
0, & k>n,
\end{cases}
\end{align*}
其中$\varepsilon _k=\begin{cases}
\frac{\overline{\alpha _k}}{|\alpha _k|},&\alpha _k\ne 0,\\
0,&\alpha _k=0\\
\end{cases}$.直接计算得
\begin{align*}
\|x^{(n)}\|_p^p=\sum\limits_{k=1}^{n}\frac{|\alpha_k|^{(q-1)p}}{\sum\limits_{j=1}^{n}|\alpha_j|^q}
=\frac{\sum\limits_{k=1}^{n}|\alpha_k|^q}{\sum\limits_{j=1}^{n}|\alpha_j|^q}=1,
\end{align*}
且
\begin{align*}
f(x^{(n)})=\sum\limits_{k=1}^{n}\alpha_k\xi_k^{(n)}
=\frac{\sum\limits_{k=1}^{n}|\alpha_k|^q}{\Bigl(\sum\limits_{j=1}^{n}|\alpha_j|^q\Bigr)^{\frac{1}{p}}}
=\Bigl(\sum\limits_{k=1}^{n}|\alpha_k|^q\Bigr)^{1-\frac{1}{p}}
=\Bigl(\sum\limits_{k=1}^{n}|\alpha_k|^q\Bigr)^{\frac{1}{q}}.
\end{align*}
因此
\begin{align*}
\parallel f\parallel =\underset{\left\| x^{\left( n \right)} \right\| =1}{\mathrm{sup}}\left| f\left( x^{\left( n \right)} \right) \right|\geqslant f(x^{(n)})=\Bigl( \sum_{k=1}^n{|\alpha _k|^q} \Bigr) ^{\frac{1}{q}}.
\end{align*}
令 $n\to\infty$ 取极限即得 $\|f\|\geqslant \|\alpha\|_q$. 结合之前已证的 $\|f\|\leqslant \|\alpha\|_q$, 最终得到 $\|f\|=\|\alpha\|_q$.
\end{enumerate}
\end{proof}

\begin{example}
如果序列$\{\alpha_k\}$使得对$\forall x=\{\xi_k\}\in l^1$，保证$\sum\limits_{k=1}^{\infty}\alpha_k\xi_k$收敛，求证：
\begin{enumerate}[(1)]
\item $\{\alpha_k\}\in l^\infty$。

\item 若$f:x\mapsto\sum\limits_{k=1}^{\infty}\alpha_k\xi_k$作为$l^1$上的线性泛函，则
\begin{align*}
\|f\|=\sup_{k\geqslant 1}|\alpha_k|.
\end{align*}
\end{enumerate}
\end{example}
\begin{proof}
\begin{enumerate}[(1)]
\item 对每个 $n\in\mathbb{N}$, 定义线性泛函 $f_n:l^1\to\mathbb{R}$ (或 $\mathbb{C}$) 为
\begin{align*}
f_n(x)=\sum\limits_{k=1}^{n}\alpha_k\xi_k,\quad x=\{\xi_k\}\in l^1.
\end{align*}
显然
\begin{align*}
|f_n(x)|\leqslant \bigl(\max_{1\leqslant k\leqslant n}|\alpha_k|\bigr)\sum\limits_{k=1}^{\infty}|\xi_k| = \bigl(\max_{1\leqslant k\leqslant n}|\alpha_k|\bigr)\|x\|_1,
\end{align*}
故 $\|f_n\|\leqslant \max_{1\leqslant k\leqslant n}|\alpha_k|$. 另一方面, 对每个 $j\leqslant n$, 取标准单位向量 $e_j\in l^1$ ,即$e_j$的第 $j$ 个分量为 $1$, 其余为 $0$, 有 $f_n(e_j)=\alpha_j$, 且 $\|e_j\|_1=1$, 因此 $\|f_n\|\geqslant |\alpha_j|$. 对所有 $1\leqslant j\leqslant n$ 取上确界得到 $\|f_n\|\geqslant \max_{1\leqslant k\leqslant n}|\alpha_k|$. 于是
\begin{align}\label{eq:fn_norm_ell1}
\|f_n\| = \max_{1\leqslant k\leqslant n}|\alpha_k|.
\end{align}
依题设, 对每个 $x\in l^1$, $\lim\limits_{n\to\infty}f_n(x)=\sum_{k=1}^{\infty}\alpha_k\xi_k$ 存在, 从而 $\{|f_n(x)|\}$ 有界. 由\hyperref[theorem:泛函分析--共鸣定理]{共鸣定理}, 存在常数 $M>0$ 使得 $\|f_n\|\leqslant M$ 对一切 $n$ 成立. 结合 \eqref{eq:fn_norm_ell1}式得
\begin{align*}
\max_{1\leqslant k\leqslant n}|\alpha_k|\leqslant M\quad \forall n\in \mathbb{N}.
\end{align*}
故 $\sup_{k\geqslant 1}|\alpha_k|\leqslant M<\infty$, 即 $\{\alpha_k\}\in l^{\infty}$.

\item 现定义 $f(x)=\sum_{k=1}^{\infty}\alpha_k\xi_k$,则
\begin{align*}
|f(x)|\leqslant \sum_{k=1}^{\infty}|\alpha_k||\xi_k|
\leqslant \bigl(\sup_{k\geqslant 1}|\alpha_k|\bigr)\sum_{k=1}^{\infty}|\xi_k|
= \bigl(\sup_{k\geqslant 1}|\alpha_k|\bigr)\|x\|_1\quad (\forall x\in l^1).
\end{align*}
所以$f$有界,即$f\in(l^1)^*$ 且 $\|f\|\leqslant \sup_{k\geqslant 1}|\alpha_k|$.
为证反向不等式, 对任意 $\varepsilon>0$, 存在 $k_0$ 使得 $|\alpha_{k_0}| > \sup_{k\geqslant 1}|\alpha_k| - \varepsilon$. 取 $x^{(0)}=e_{k_0}$, 则 $\|x^{(0)}\|_1=1$, 且 $f(x^{(0)})=\alpha_{k_0}$, 从而
\begin{align*}
\|f\|\geqslant |f(x^{(0)})| = |\alpha_{k_0}| > \sup_{k\geqslant 1}|\alpha_k| - \varepsilon.
\end{align*}
由 $\varepsilon$ 的任意性得 $\|f\|\geqslant \sup_{k\geqslant 1}|\alpha_k|$. 综上所述,
\begin{align*}
\|f\| = \sup_{k\geqslant 1}|\alpha_k|.
\end{align*}
\end{enumerate}
\end{proof}

\begin{example}
设$\mathscr{X},\mathscr{Y}$是$B$空间，$A\in\mathscr{L}(\mathscr{X},\mathscr{Y})$是满射的。求证：如果在$\mathscr{Y}$中$y_n\to y_0$，则$\exists C>0$与$x_n\to x_0$，使得$Ax_n=y_n$，且$\|x_n\|\leqslant C\|y_n\|$。
\end{example}
\begin{proof}

由\cref{proposition:开映射--有界原像}知
\begin{align}\label{eq:existence}
\forall y\in\mathscr{Y},\ \exists x\in\mathscr{X}\text{ 满足 } Ax=y \text{ 且 } \|x\|\leqslant M\|y\|.
\end{align}
现设 $y_n\to y_0$。由 \eqref{eq:existence}式，存在 $x_0\in\mathscr{X}$ 使得 $Ax_0=y_0$ 且 $\|x_0\|\leqslant M\|y_0\|$。
对每个 $n$，对向量 $y_n-y_0$ 应用 \eqref{eq:existence}式，存在 $u_n\in\mathscr{X}$ 满足
\[
Au_n=y_n-y_0,\qquad \|u_n\|\leqslant M\|y_n-y_0\|.
\]
令$x_n\triangleq x_0+u_n,$则 $Ax_n=Ax_0+Au_n=y_0+(y_n-y_0)=y_n$，并且
\[
\|x_n-x_0\|=\|u_n\|\leqslant M\|y_n-y_0\|\to0,
\]
故 $x_n\to x_0$。
下面证明存在常数 $C>0$ 使得 $\|x_n\|\leqslant C\|y_n\|$ 对一切 $n$ 成立。分两种情形：

情形 1：$y_0\neq0$。
由 $y_n\to y_0$，存在 $N$，当 $n>N$ 时，有
\[
\|y_n-y_0\|\leqslant \|y_0\|,\qquad \|y_n\|\geqslant \frac12\|y_0\|.
\]
于是当 $n>N$ 时，有
\[
\|x_n\|\leqslant \|x_0\|+\|u_n\|
\leqslant M\|y_0\|+M\|y_n-y_0\|
\leqslant M\|y_0\|+M\|y_0\|=2M\|y_0\|
\leqslant 2M\cdot 2\|y_n\|=4M\|y_n\|.
\]
若$y_n=0(1\leqslant n\leqslant N)$，则令$C_1=0.$否则,令$C_1\triangleq \max_{\substack{1\leqslant n\leqslant N\\ y_n\neq0}}\frac{\|x_n\|}{\|y_n\|}.$
再取
\begin{align*}
C\triangleq \max\{4M,\,C_1\},
\end{align*}
则对所有 $n$ 均有 $\|x_n\|\leqslant C\|y_n\|$。

情形 2：$y_0=0$。
此时由 \eqref{eq:existence} 式,可直接选取$C=M$,则对所有 $n$ 均有 $\|x_n\|\leqslant C\|y_n\|$。

综上所述，总存在常数 $C>0$ 及满足 $Ax_n=y_n$ 的序列 $x_n\to x_0$，使得 $\|x_n\|\leqslant C\|y_n\|$ 对所有 $n$ 成立.
\end{proof}

\begin{example}
设$\mathscr{X},\mathscr{Y}$是$B$空间，$T$是闭线性算子，$D(T)\subseteq\mathscr{X}$，$R(T)\subseteq\mathscr{Y}$，$N(T)\triangleq\{x\in\mathscr{X}\mid Tx=\theta\}$。
\begin{enumerate}[(1)]
\item 求证：$N(T)$是$\mathscr{X}$的闭线性子空间。
\item 求证：$N(T)=\{\theta\},R(T)$在$\mathscr{Y}$中闭的充要条件是，$\exists\alpha>0$，使得
\begin{align*}
\|x\|\leqslant\alpha\|Tx\|\quad (\forall x\in D(T)).
\end{align*}
\item 如果用$d(x,N(T))$表示点$x\in\mathscr{X}$到集合$N(T)$的距离$\left(\inf\limits_{z\in N(T)}\|z-x\|\right)$。求证：$R(T)$在$\mathscr{Y}$中闭的充要条件是，$\exists\alpha>0$，使得
\begin{align*}
d(x,N(T))\leqslant\alpha\|Tx\|\quad (\forall x\in D(T)).
\end{align*}
\end{enumerate}
\end{example}
\begin{proof}
\begin{enumerate}[(1)]
\item 显然 $N(T)$ 是线性子空间. 下证其闭性. 设 $\{x_n\}\subseteq N(T)$ 且 $x_n\to x\in\mathscr{X}$. 由 $x_n\in D(T)$, $Tx_n=\theta \to\theta$, 而 $T$ 是闭算子, 故 $x\in D(T)$ 且 $Tx=\theta$, 即 $x\in N(T)$. 所以 $N(T)$ 是闭线性子空间.

\item {\heiti 必要性:} 设 $N(T)=\{\theta\}$ 且 $R(T)$ 闭,$T$是单射,从而$T$ 是 $D(T)$ 到 $R(T)$ 的双射. 令 $T^{-1}:R(T)\to D(T)$ 为其逆映射. 因 $T$ 是闭算子,故由\cref{proposition:有界线性算子必是闭的}知$T^{-1}$ 也是闭算子. 而 $R(T)\subseteq\mathscr{Y}$ 是闭集, $\mathscr{Y}$ 完备, 故 $R(T)$ 是 Banach 空间. 由\hyperref[theorem:泛函分析--闭图像定理]{闭图像定理}, $T^{-1}$ 连续, 即存在 $\alpha>0$ 使得对任意 $y\in R(T)$, $\|T^{-1}y\|\leqslant\alpha\|y\|$. 取 $x=T^{-1}y$, 则 $y=Tx$, 得 $\|x\|\leqslant\alpha\|Tx\|$ 对一切 $x\in D(T)$ 成立.

{\heiti 充分性:} 假设存在 $\alpha>0$ 使得 $\|x\|\leqslant\alpha\|Tx\|$ 对一切 $x\in D(T)$ 成立. 若 $Tx=\theta$, 则 $\|x\|=0$, 故 $x=\theta$, 即 $N(T)=\{\theta\}$. 再证 $R(T)$ 闭. 设 $\{y_n\}\subseteq R(T)$ 且 $y_n\to y_0\in\mathscr{Y}$. 取 $x_n\in D(T)$ 使得 $Tx_n=y_n$. 由假设得
\begin{align*}
\|x_n-x_m\|\leqslant\alpha\|Tx_n-Tx_m\|=\alpha\|y_n-y_m\|,
\end{align*}
故 $\{x_n\}$ 是 $\mathscr{X}$ 中的 Cauchy 列.由 $\mathscr{X}$ 完备, 故存在 $x_0\in\mathscr{X}$ 使得 $x_n\to x_0$. 此时 $Tx_n=y_n\to y_0$, 又$T$ 是闭算子, 故 $x_0\in D(T)$ 且 $Tx_0=y_0$, 于是 $y_0\in R(T)$. 所以 $R(T)$ 闭.

\item 考虑商空间 $\widetilde{\mathscr{X}}=\mathscr{X}/N(T)$. 因 $N(T)$ 闭,所以由\cref{theorem:闭线性子空间的商空间完备}知$\widetilde{\mathscr{X}}$ 在商范数 $\|[x]\|=d(x,N(T))$ 下是 Banach 空间. 定义算子 $\widetilde{T}:D(\widetilde{T})\subseteq\widetilde{\mathscr{X}}\to\mathscr{Y}$ 如下:
\begin{align*}
D(\widetilde{T})=\{[x]: x\in D(T)\},\quad \widetilde{T}([x])=Tx.
\end{align*}
若 $x_1,x_2\in D(T)$ 满足 $[x_1]=[x_2]$, 则 $x_1-x_2\in N(T)$, 从而 $Tx_1=Tx_2$, 故 $\widetilde{T}$ 良定义且线性. 现证 $\widetilde{T}$ 是闭算子. 设 $\{[x_n]\}\subseteq D(\widetilde{T})$ 满足 $[x_n]\to [x]\in\widetilde{\mathscr{X}}$ 且 $\widetilde{T}([x_n])\to y\in\mathscr{Y}$.取定代表元 $x\in [x]$, 由商范数定义知
\begin{align*}
\underset{z\in N\left( T \right)}{\mathrm{inf}}\left\| \left( x_n-x \right) -z \right\| =d(x_n-x,N(T))=\|[x_n]-[x]\|\to0.
\end{align*}
于是对每个 $n$, 可选取 $z_n\in N(T)$ 使得 $\|(x_n-x)-z_n\|<\frac{1}{n}+\|[x_n]-[x]\|$, 从而 $\|(x_n-z_n)-x\|\to0$. 令 $\widetilde{x}_n=x_n-z_n$, 则 $[\widetilde{x}_n]=[x_n]$, 且 $\widetilde{x}_n\to x$. 同时 $T\widetilde{x}_n=Tx_n=\widetilde{T}([x_n])\to y$. 由于 $T$ 是闭算子, 得 $x\in D(T)$ 且 $Tx=y$.  因此 $[x]\in D(\widetilde{T})$ 且 $\widetilde{T}([x])=y$, 故 $\widetilde{T}$是闭的. 

若 $\widetilde{T}([x])=\theta$, 则 $Tx=\theta$, 即 $x\in N(T)$, $[x]=[0]$.故$N(\widetilde{T})=\{\theta\}$.

{\heiti 必要性:} 若 $R(T)$ 闭, 则 $R(\widetilde{T})=R(T)$ 是 Banach 空间. 由 (2) 的结论知存在 $\alpha>0$ 使得
\begin{align*}
d(x,N(T))=\|[x]\|\leqslant\alpha\|\widetilde{T}([x])\|\quad(\forall x\in D(T))=\alpha\|Tx\|.
\end{align*}

{\heiti 充分性:} 若存在 $\alpha>0$ 使 $d(x,N(T))\leqslant\alpha\|Tx\|$, 则 $\|[x]\|\leqslant\alpha\|\widetilde{T}([x])\|$. 由 (2)的结论知 $R(\widetilde{T})$ 在 $\mathscr{Y}$ 中闭, 即 $R(T)$ 闭.
\end{enumerate}
\end{proof}

\begin{example}
设$H$是Hilbert空间，$A\in \mathscr{L}(H)$，并且存在$m>0$，使得
\begin{align*}
\left|\left(Ax,x\right)\right|\geqslant m\left\|x\right\|^2\quad(\forall x\in H).
\end{align*}
其中$\mathscr{L}(H)$表示$H$到$H$的有界线性算子全体。证明：存在唯一的$ A^{-1}\in \mathscr{L}(H)$,且满足
\begin{align*}
\|B^{-1}\|\leqslant \frac{1}{m}.
\end{align*}
\end{example}
\begin{proof}
定义
\begin{align*}
a(x,y) \triangleq (x,Ay) \quad (\forall x,y\in H),
\end{align*}
则显然$a(x,y)$是$H$上的共轭双线性函数.
下验证$a(x,y)$满足\hyperref[theorem:泛函分析--Lax-Milgram定理]{Lax-Milgram定理}的两个条件.
首先, 由Cauchy-Schwarz不等式及$A$的有界性, 有
\begin{align*}
|a(x,y)| = |(x,Ay)| \leqslant \|x\|\|Ay\| \leqslant \|A\|\|x\|\|y\|,
\end{align*}
其次, 由题设可得
\begin{align*}
|a(x,x)| = |(x,Ax)| = |(Ax,x)| \geqslant m\|x\|^2,
\end{align*}
于是由\hyperref[theorem:泛函分析--Lax-Milgram定理]{Lax-Milgram定理}, 存在唯一的有连续逆的连续线性算子$B\in\mathscr{L}(H)$, 满足
\begin{align*}
\|B^{-1}\|\leqslant \frac{1}{m},\quad a(x,y) = (x,By) \quad (\forall x,y\in H).
\end{align*}
而$a(x,y) = (x,Ay)$, 故
\begin{align*}
(x,By) = (x,Ay) \quad (\forall x,y\in H).
\end{align*}
对$\forall y\in H$,取$x=By-Ay$,得$(By-Ay,By-Ay)=0$,进而$By = Ay$对任意$y\in H$成立, 即$B = A$. 因此$A$存在唯一的连续逆$A^{-1} = B^{-1}\in\mathscr{L}(H)$,且满足
\begin{align*}
\|B^{-1}\|\leqslant \frac{1}{m}.
\end{align*}
\end{proof}

\begin{proposition}\label{proposition:共轭双线性型的Riesz表示定理}
设$a(x,y)$是Hilbert空间$H$上的一个共轭双线性泛函，满足：
\begin{enumerate}[(1)]
\item $\exists M>0$，使得$|a(x,y)|\leqslant M\|x\|\cdot\|y\|\quad (\forall x,y\in H)$;
\item $\exists\delta>0$，使得$|a(x,x)|\geqslant\delta\|x\|^2\quad (\forall x\in H)$.
\end{enumerate}
求证：$\forall f\in H^*,$存在唯一的$y_f\in H$，使得
\begin{align*}
a(x,y_f)=f(x)\quad (\forall x\in H),
\end{align*}
而且$y_f$连续地依赖于$f$,并且映射$f\mapsto y_f$是Lipschitz连续的.
\end{proposition}
\begin{proof}

固定$y\in H$，由条件(1)知，对任意$x\in H$有$|a(x,y)|\leqslant M\|x\|\|y\|$，从而$x\mapsto a(x,y)$是$H$上的有界线性泛函. 根据Riesz表示定理，存在唯一的$a_y\in H$，使得
\begin{align*}
a(x,y)=\langle x, a_y\rangle \quad (\forall x\in H).
\end{align*}
定义映射
\begin{align*}
A:H\to H,\quad Ay= a_y.
\end{align*}
先证$A$是线性算子. 对任意$\alpha,\beta\in\mathbb{C}$及$y_1,y_2\in H$，有
\begin{align*}
\langle x, A(\alpha y_1+\beta y_2)\rangle &= a(x,\alpha y_1+\beta y_2) = \overline{\alpha}a(x,y_1)+\overline{\beta}a(x,y_2) \\
&= \overline{\alpha}\langle x, Ay_1\rangle+\overline{\beta}\langle x, Ay_2\rangle = \langle x, \alpha Ay_1+\beta Ay_2\rangle \quad (\forall x\in H),
\end{align*}
由唯一性得$A(\alpha y_1+\beta y_2)=\alpha Ay_1+\beta Ay_2$，故$A$线性.

其次证$A$有界. 对任意$y\in H$，
\begin{align*}
\|Ay\| = \sup_{\|x\|=1}|\langle x, Ay\rangle| = \sup_{\|x\|=1}|a(x,y)| \leqslant \sup_{\|x\|=1} M\|x\|\|y\| = M\|y\|,
\end{align*}
所以$\|A\|\leqslant M$，也即$A$连续.

再由条件(2)推出下界估计. 对任意$y\in H$，$a(y,y)=\langle y, Ay\rangle$，再由\hyperref[theorem:泛函分析--Cauchy-Schwarz不等式]{Cauchy-Schwarz不等式}可得
\begin{align*}
\delta\|y\|^2 \leqslant |a(y,y)| = |\langle y, Ay\rangle| \leqslant \|y\|\|Ay\|,
\end{align*}
当$y\neq 0$时约去$\|y\|$得$\|Ay\|\geqslant \delta\|y\|$；$y=0$时显然$\|Ay\|\geqslant \delta\|y\|$也成立. 因此
\begin{align}\label{eq:okokkooooo}
\|Ay\|\geqslant \delta\|y\|,\quad \forall y\in H.
\end{align}
若$Ay=0$则$\|y\|=0$，即$y=0$，从而$A$是单射。若$\{Ay_n\}\subseteq R(A)$收敛于某$z\in H$，则$\{Ay_n\}$是Cauchy列。由\eqref{eq:okokkooooo}式知
\begin{align*}
\|y_n-y_m\|\leqslant\frac{1}{\delta}\|Ay_n-Ay_m\|.
\end{align*}
进而由$\{Ay_n\}$是Cauchy列知$\{y_n\}$也是Cauchy列，设$y_n\to y$，由$A$的连续性得$Ay_n\to Ay$，故$z=Ay\in R(A)$.因此值域$R(A)$是闭集.

下证$A$是满射. 设$z\in(R(A))^\perp$，即对任意$y\in H$有$\langle Ay, z\rangle =0$. 取$y=z$得$\langle Az, z\rangle =0$. 另一方面，
\begin{align*}
|\langle Az, z\rangle| = |\overline{\langle z, Az\rangle}| = |\langle z, Az\rangle| = |a(z,z)| \geqslant \delta\|z\|^2.
\end{align*}
联合两式得$0\geqslant \delta\|z\|^2$，故$z=0$. 这证明了$R(A)$的正交补为零，因此由\subcref{proposition:Hilbert空间稠密线性空间正交补为0}{proposition:Hilbert空间稠密线性空间正交补为0-2}知$R(A)$在$H$中稠密. 结合$R(A)$为闭，即得$R(A)=H$，即$A$满射. 于是$A$是$H$到$H$的双射，且由\eqref{eq:okokkooooo}式知逆算子$A^{-1}$存在且有界，$\|A^{-1}\|\leqslant\frac{1}{\delta}$.

现在对任意$f\in H^*$，由\hyperref[theorem:Riesz表示定理(Hilbert空间)]{Riesz表示定理}，存在唯一的$z_f\in H$使得
\begin{align}\label{eq------32393333}
\|z_f\|=\|f\|_{H^*},\quad f(x)=\langle x, z_f\rangle \quad (\forall x\in H).
\end{align}
令$y_f = A^{-1}z_f$，则对任意$x\in H$，
\begin{align*}
a(x,y_f)=\langle x, Ay_f\rangle = \langle x, z_f\rangle = f(x).
\end{align*}
若另有$\widetilde{y}\in H$满足$a(x,\widetilde{y})=f(x)$对一切$x$成立，则相减得$a(x,\widetilde{y}-y_f)=0$对所有$x$成立. 取$x=\widetilde{y}-y_f$，利用条件(2)得
\begin{align*}
\delta\|\widetilde{y}-y_f\|^2 \leqslant |a(\widetilde{y}-y_f,\widetilde{y}-y_f)| = 0,
\end{align*}
所以$\widetilde{y}=y_f$，这证明了$y_f$的唯一性.

最后证$y_f$连续依赖于$f$. 由\eqref{eq------32393333}式知$\|z_f\|=\|f\|_{H^*}$，故
\begin{align*}
\|y_f\| = \|A^{-1}z_f\| \leqslant \|A^{-1}\|\,\|z_f\| \leqslant\frac{1}{\delta}\|f\|_{H^*}.
\end{align*}
因此映射$f\mapsto y_f$是Lipschitz连续的.
\end{proof}

\begin{example}
设$\Omega$是$\mathbb{R}^2$中边界光滑的有界开区域，$\alpha:\Omega\to\mathbb{R}$有界可测并满足$0<\alpha_0\leqslant\alpha,f\in L^2(\Omega)$。规定：
\begin{gather*}
a(u,v)\triangleq\int_{\Omega}(\nabla u\cdot\nabla v+\alpha uv)\mathrm{d}x\mathrm{d}y\quad (\forall u,v\in H^1(\Omega)),\\
F(v)\triangleq\int_{\Omega}f\cdot v\mathrm{d}x\mathrm{d}y\quad (\forall v\in L^2(\Omega)).
\end{gather*}
求证：存在唯一的$u\in H^1(\Omega)$满足
\begin{align*}
a(u,v)=F(v)\quad (\forall v\in H^1(\Omega)).
\end{align*}
\end{example}
\begin{proof}
由\cref{theorem:Sobolev空间的完备性}知,$H^1(\Omega)$按通常内积构成Hilbert空间. 由题设，$a(\cdot,\cdot)$是$H^1(\Omega)$上的共轭双线性泛函. 对任意$u,v\in H^1(\Omega)$，
\begin{align*}
|a(u,v)| &\leqslant \int_{\Omega}|\nabla u\cdot\nabla v|\mathrm{d}x\mathrm{d}y+\int_{\Omega}|\alpha uv|\mathrm{d}x\mathrm{d}y\\
&\leqslant \|\nabla u\|_{L^2}\|\nabla v\|_{L^2}+\|\alpha\|_{L^\infty}\|u\|_{L^2}\|v\|_{L^2}\\
&\leqslant \max\{1,\|\alpha\|_{L^\infty}\}\|u\|_{H^1}\|v\|_{H^1},
\end{align*}
故存在$M=\max\{1,\|\alpha\|_{L^\infty}\}$使得$|a(u,v)|\leqslant M\|u\|\|v\|$. 又由$0<\alpha_0\leqslant\alpha$知
\begin{align*}
a(u,u)=\int_{\Omega}(|\nabla u|^2+\alpha u^2)\mathrm{d}x\mathrm{d}y\geqslant \int_{\Omega}(|\nabla u|^2+\alpha_0 u^2)\mathrm{d}x\mathrm{d}y\geqslant \min\{1,\alpha_0\}\|u\|_{H^1}^2,
\end{align*}
取$\delta=\min\{1,\alpha_0\}$即得$|a(u,u)|\geqslant\delta\|u\|^2$. 因此$a$满足\cref{proposition:共轭双线性型的Riesz表示定理}中的全部条件. 再证$F\in (H^1(\Omega))^*$. 由$f\in L^2(\Omega)$及\hyperref[theorem:泛函分析--Cauchy-Schwarz不等式]{Cauchy-Schwarz不等式}得
\begin{align*}
|F(v)|\leqslant \|f\|_{L^2}\|v\|_{L^2}\leqslant \|f\|_{L^2}\|v\|_{H^1}\quad (\forall v\in H^1(\Omega)),
\end{align*}
所以$F$是$H^1(\Omega)$上的有界线性泛函. 应用\cref{proposition:共轭双线性型的Riesz表示定理}的结论，存在唯一的$u\in H^1(\Omega)$使得
\begin{align*}
a(u,v)=F(v)\quad (\forall v\in H^1(\Omega)).
\end{align*}
\end{proof}




\end{document}